% explainer-source-hash: f853827457ce4135f852201be81a6e16bb0e75ff89457f8e3540348443911a82
% explainer-source-commit: 63f18f95db27c45910ef56ad52463f8853a5e12a
% explainer-generated: 2026-07-16
% Regenerate with /explain-all (see WIRING.md). Do not edit this stamp by hand:
% it is how the toolchain knows whether this document still matches the code.
\documentclass[11pt,a4paper]{article}
\input{preamble}

\begin{document}

\begin{titlepage}
\centering
\vspace*{2.5cm}
{\Huge\bfseries\color{inkblue} urdu-slm\par}
\vspace{0.6cm}
{\LARGE\color{accent} Deep Dive\par}
\vspace{1.2cm}
{\large A small Urdu language model trained from scratch:\\
the whole stack, from raw public text to a decoder-only transformer\par}
\vspace{2cm}
\begin{minipage}{0.82\textwidth}
\small
This document explains every part of the \file{urdu-slm} repository to a reader who
has never trained a language model. It defines each technical term the first time it
appears, prefers a diagram to a paragraph wherever a diagram can carry the idea, and
separates what was verified by running the code from what was read in the README.

Where the README and the code disagree, the code wins, and the disagreement is
recorded rather than smoothed over. Section~\ref{sec:honest} collects every such
finding in one place.
\end{minipage}
\vfill
{\small Working document for the author's own understanding and interview defense.\par}
\end{titlepage}

\tableofcontents
\newpage

% =====================================================================
\section{Orientation: what this repository actually is}

\file{urdu-slm} builds a language model for Urdu starting from nothing but publicly
downloadable text. It is not a fine-tune of somebody else's model, and it does not
import modeling code from the \code{transformers} library. Every layer of the network
is written out in PyTorch in one 186-line file, and every stage that turns a raw
Wikipedia dump into training-ready numbers is written out in the \file{pipeline/}
package.

\begin{keyidea}{The whole project in five moves}
\begin{enumerate}
\item \textbf{Collect} raw Urdu text from public sources that need no login.
\item \textbf{Clean} it: normalize the spelling, throw away non-Urdu and boilerplate,
      and remove duplicates.
\item \textbf{Build a tokenizer}, a scheme for cutting Urdu text into reusable chunks,
      trained on that cleaned Urdu rather than borrowed from an English model.
\item \textbf{Train a transformer} to predict the next chunk, over and over, on a rented GPU.
\item \textbf{Measure} it honestly, and report the numbers that came out, including the
      unflattering ones.
\end{enumerate}
\end{keyidea}

The repository ships three trained artifacts: a 13M-parameter model trained on a CPU as
a proof that the pipeline works end to end, and two 124M-parameter models trained on a
rented A100 GPU (\code{base} on 2026-07-07 and \code{base\_v2} on 2026-07-08). The
per-step training logs for all three are committed, which is what made much of this
document checkable rather than merely repeated.

\subsection{A necessary note about the script}
\label{sec:script}

\begin{honest}{Why you will not see a single Urdu character in this PDF}
Urdu is written in the Perso-Arabic script: right to left, with letters that change
shape depending on their neighbours. This document is typeset with \code{pdflatex},
which on this machine has no font capable of rendering that script and no engine for
right-to-left shaping. Pasting Urdu text into a code listing here is not a cosmetic
problem, it is a hard build failure: the listings package reads raw bytes and the input
encoder aborts on the first multi-byte character.

So every Urdu string in this document is handled one of three ways: transliterated into
Latin letters, described in English, or given as Unicode codepoints. The real script
lives in the repository's own README and in \file{evaldata/cloze.jsonl}, which render
correctly in any modern editor. Nothing has been hidden; it simply cannot be drawn here,
and an honest sentence beats a document that will not build.
\end{honest}

% =====================================================================
\section{Part I: Foundations, assuming nothing}

If you already know what a decoder-only transformer is, skip to
Section~\ref{sec:glance}. Nothing in this part is specific to this repository; it is the
vocabulary the rest of the document uses.

\subsection{What a language model is}

\begin{jargon}{Language model}
A program that, given a stretch of text, assigns a probability to every possible next
piece of text. That is the entire definition. It does not "understand" in any agreed
sense; it estimates how likely each continuation is, given everything it has read
during training.
\end{jargon}

Suppose the text so far is "the capital of Pakistan is". A language model outputs a
number for every entry in its vocabulary: perhaps $0.6$ for "Islamabad", $0.05$ for
"Karachi", and a vanishingly small number for "bicycle". Those numbers must sum to one
across the whole vocabulary. To generate text you sample from that distribution, append
what you drew, and repeat.

\begin{jargon}{Training}
The process of adjusting the model's internal numbers (its \emph{parameters}) so that
the probabilities it assigns to real text are as high as possible. You show it real
sentences with the answer hidden, measure how surprised it was by the true next piece,
and nudge every parameter slightly in the direction that would have reduced that
surprise.
\end{jargon}

\begin{jargon}{Parameter}
One adjustable number inside the model. The \code{base} model here has 123.7 million of
them. They start as random noise and end, after training, holding everything the model
knows. "124M parameters" is a statement about the model's capacity to store patterns,
not about its size on disk (though the two are related: 123.7M parameters at 4 bytes
each is the 495MB of \file{runs/base\_v2/model.pt}).
\end{jargon}

\begin{jargon}{Loss, and cross-entropy}
The number that measures "how surprised the model was". This project uses
\emph{cross-entropy loss}, measured in \emph{nats}. If the model gave the true next
token probability $p$, the loss for that token is $-\ln p$. A confident correct
prediction ($p$ near 1) gives a loss near 0; assigning the truth a probability of
one-in-a-thousand gives a loss of about 6.9. Training is the search for parameters that
make the average loss small.
\end{jargon}

\begin{jargon}{Perplexity}
Loss made human-readable: $\mathrm{perplexity} = e^{\mathrm{loss}}$. It answers "on average,
how many equally-likely options did the model feel it was choosing between?" A
perplexity of 28.93 (this project's best model) means the model was about as uncertain
as someone picking uniformly among 29 choices. Lower is better. A perplexity equal to
the vocabulary size (32,000) is a model that has learned nothing.
\end{jargon}

\subsection{Tokens: why models do not read letters}

A model cannot consume text directly; it consumes integers. The mapping between the two
is where this project does its most distinctive work.

\begin{jargon}{Token}
The unit a language model reads and writes. Not a character, and usually not a word: a
frequently occurring \emph{chunk} of characters. In an English model, "the" is typically
one token, and an unusual word like "tokenization" might be three ("token", "iz",
"ation"). The model's vocabulary is the fixed list of every chunk it knows, here 32,000
of them.
\end{jargon}

\begin{jargon}{Tokenizer}
The program that cuts text into tokens and back again. It is trained, separately from
and before the model, by looking at a corpus and deciding which chunks are worth having
their own vocabulary entry. Once trained it is frozen: the model is built around one
specific tokenizer and cannot be used with another.
\end{jargon}

Why does chunking matter so much? Because every cost in the system is counted in tokens.
The GPU bill is per token. The model's context window, how much text it can see at once,
is measured in tokens (1024 here). If a tokenizer needs five times as many tokens to
express the same sentence, then the same GPU budget reads one fifth as much text, and
the model's 1024-token window holds one fifth as much meaning.

\begin{keyidea}{Byte-Pair Encoding (BPE), the algorithm used here}
Start with a vocabulary of single characters (or, here, single bytes). Then repeat: scan
the corpus, find the pair of adjacent symbols that occurs most often, and merge it into
one new symbol added to the vocabulary. Do this until the vocabulary reaches the target
size (32,000). Common sequences in the training text end up as single tokens; rare ones
stay split. The tokenizer is therefore a direct fingerprint of the language it was
trained on.
\end{keyidea}

\subsection{Why tokenizing Urdu is not like tokenizing English}
\label{sec:whyurdu}

This is the heart of the project's motivation, so it is worth being precise.

Text is stored as Unicode codepoints, encoded into bytes with UTF-8. UTF-8 is
deliberately biased toward English: every ASCII character (a to z, digits, punctuation)
costs exactly \textbf{one byte}. Urdu's letters live in the Arabic block, codepoints
U+0600 to U+06FF, and every one of them costs \textbf{two bytes}.

Now consider a byte-level BPE tokenizer trained on English web text, which is what
GPT-2's is. Its 50,257 merges were chosen to compress English. It has a token for "the"
and a token for " Pakistan". It has essentially no merges covering Urdu byte sequences,
because it barely saw any. So when it meets Urdu it falls back to its base alphabet and
emits roughly one token per byte, and each Urdu letter is two bytes.

\begin{keyidea}{The compounding penalty}
An English tokenizer on Urdu pays twice over. First, Urdu characters cost two UTF-8
bytes each instead of one. Second, no merges fire, so those bytes never combine back
into larger chunks. The result is a tokenizer that spends more than two tokens per Urdu
\emph{character}, while a tokenizer trained on Urdu spends about one token per four
characters. The measured gap is Section~\ref{sec:tokclaim}, and it is 5.23x.
\end{keyidea}

Two further properties of Urdu shape the pipeline in Section~\ref{sec:pipeline}:

\begin{itemize}
\item \textbf{Encoding variants.} The same visible Urdu letter can be written with
      different codepoints, because Urdu borrows the Arabic script and both the Arabic
      and Persian-adapted forms circulate. The letter "kaf" appears as U+0643 (Arabic)
      or U+06A9 (Urdu keheh); "yeh" as U+064A (Arabic) or U+06CC (Farsi). To a human
      these look identical. To a BPE trainer they are unrelated symbols, so it would
      waste vocabulary learning both spellings of every word. \file{normalize.py} folds
      them; see Section~\ref{sec:normalize}.
\item \textbf{Rich morphology.} Urdu inflects heavily, so a word-level vocabulary would
      explode. Subword BPE handles this naturally by learning stems and endings.
\end{itemize}

\subsection{What a transformer is}

The model itself is a \emph{decoder-only transformer}. Taking that apart:

\begin{jargon}{Neural network}
A stack of layers, each of which multiplies its input by a matrix of parameters and
applies a simple non-linear function. Numbers go in at the bottom, get reshaped layer by
layer, and predictions come out at the top. Training adjusts the matrices.
\end{jargon}

\begin{jargon}{Embedding}
The lookup table that turns a token id (an integer) into a vector of numbers the network
can work with. Here it maps each of 32,000 token ids to a vector of 768 numbers for the
\code{base} model. Those 32{,}000 $\times$ 768 = 24.6 million numbers are themselves
parameters, learned during training. This table's size is a recurring theme in
Section~\ref{sec:presets}.
\end{jargon}

\begin{jargon}{Attention}
The mechanism that lets a position in the sequence look at other positions. For each
position the model computes three vectors: a \emph{query} (what am I looking for?), a
\emph{key} (what do I offer?), and a \emph{value} (what do I pass along if chosen?).
Each position compares its query against every key, turns the comparison scores into
weights that sum to one, and takes that weighted average of the values. This is how the
word "it" finds the noun it refers to twelve words back.
\end{jargon}

\begin{jargon}{Causal (masked) attention, and "decoder-only"}
"Causal" means a position may attend only to itself and to positions \emph{before} it,
never after. This is what makes next-token prediction a fair game: when predicting
position 50, the model must not be allowed to peek at position 51, which is the answer.
"Decoder-only" describes an architecture built entirely from causal blocks, with no
separate encoder. Every modern text-generating model (the GPT family, Llama, and this
one) is decoder-only.
\end{jargon}

Figure~\ref{fig:causal} shows what the mask actually does. This property is not merely
asserted in the repository; \file{tests/test\_model.py} pins it with a direct
experiment, described in Section~\ref{sec:tests}.

\begin{figure}[H]
\centering
\resizebox{0.86\textwidth}{!}{
\begin{tikzpicture}[x=1.0cm,y=1.0cm]
  \foreach \i in {0,1,2,3,4} {
    \foreach \j in {0,1,2,3,4} {
      \pgfmathparse{int(\j<=\i)}
      \ifnum\pgfmathresult=1
        \fill[accent!35] (\j,-\i) rectangle ++(0.9,-0.9);
        \draw[inkblue,line width=0.7pt,line cap=round,line join=round]
          (\j+0.28,-\i-0.46) -- (\j+0.41,-\i-0.60) -- (\j+0.64,-\i-0.30);
      \else
        \fill[softgrey] (\j,-\i) rectangle ++(0.9,-0.9);
        \draw[warnred!70,line width=0.6pt,line cap=round]
          (\j+0.31,-\i-0.31) -- (\j+0.59,-\i-0.59)
          (\j+0.59,-\i-0.31) -- (\j+0.31,-\i-0.59);
      \fi
      \draw[linegrey] (\j,-\i) rectangle ++(0.9,-0.9);
    }
  }
  \foreach \j/\t in {0/t1,1/t2,2/t3,3/t4,4/t5}
    \node[font=\scriptsize\ttfamily,text=inkblue] at (\j+0.45,0.35) {\t};
  \foreach \i/\t in {0/t1,1/t2,2/t3,3/t4,4/t5}
    \node[font=\scriptsize\ttfamily,text=inkblue,anchor=east] at (-0.15,-\i-0.45) {\t};
  \node[font=\small\itshape,text=black!65] at (2.25,1.0) {key positions (what can be looked at)};
  \node[font=\small\itshape,text=black!65,rotate=90] at (-1.2,-2.25) {query positions};
  \node[box,text width=5.2cm,anchor=west] at (6.0,-1.0)
    {Position \code{t3} may attend to\\ \code{t1}, \code{t2}, \code{t3} only.};
  \node[extbox,text width=5.2cm,anchor=west] at (6.0,-3.2)
    {The greyed cells are forced to zero\\ weight. Without them the model\\ could read the answer it is\\ being asked to predict.};
\end{tikzpicture}}
\caption{Causal masking. Each row is a position asking a question; each column is a
position that could answer. Everything above the diagonal is forbidden.}
\label{fig:causal}
\end{figure}

\begin{jargon}{Context window (block size)}
The maximum number of tokens the model can attend over at once, \code{max\_seq\_len} in
the code: 512 for \code{tiny}, 1024 for the others. Attention cost grows with the square
of this number, which is why it is a deliberate, budgeted choice rather than "as large as
possible".
\end{jargon}

\newpage
% =====================================================================
\section{Part II: The system at a glance}
\label{sec:glance}

\begin{figure}[H]
\centering
\resizebox{\textwidth}{!}{
\begin{tikzpicture}[node distance=7mm]
  \node[extbox,text width=3.0cm] (wiki) {Urdu Wikipedia\\ XML dump (bz2)\\ \scriptsize CC BY-SA 4.0};
  \node[extbox,text width=3.0cm,below=of wiki] (leip) {Leipzig newscrawl\\ 2016 + 2011 (tar.gz)\\ \scriptsize CC BY-NC 4.0};
  \node[extbox,text width=3.0cm,below=of leip] (cc) {CC-100 Urdu\\ ur.txt.xz \scriptsize (v2 only)};
  \node[extbox,text width=3.0cm,below=of cc] (wd) {Wikidata SPARQL\\ \scriptsize CC0, planned run};
  \node[lbl,below=3mm of wd,text width=3.4cm] (leg) {\scriptsize dashed = built and tested,\\ not yet run};

  \node[box,right=14mm of leip,text width=2.9cm] (src) {\code{sources.py}\\ \scriptsize one reader per format,\\ streams documents};
  \node[box,right=8mm of src,text width=2.9cm] (norm) {\code{normalize.py}\\ \scriptsize NFC, fold variants,\\ script-ratio filter};
  \node[box,right=8mm of norm,text width=2.9cm] (dd) {\code{dedup.py}\\ \scriptsize blake2b exact +\\ MinHash/LSH near};
  \node[box,right=8mm of dd,text width=2.7cm] (sp) {split\\ \scriptsize deterministic\\ 99/1 train/val};

  \node[store,below=12mm of sp,text width=3.0cm] (jsonl) {\file{data/\{train,val\}/}\\ \scriptsize *.jsonl shards};

  \node[box,below=12mm of jsonl,text width=3.1cm] (tt) {\code{train\_tokenizer.py}\\ \scriptsize 32k byte-level BPE};
  \node[store,left=10mm of tt,text width=2.6cm] (tokj) {\file{tokenizer/}\\ \file{tokenizer.json}};
  \node[box,left=10mm of tokj,text width=3.1cm] (tc) {\code{tokenize\_corpus.py}\\ \scriptsize encode, upsample,\\ filter by source};
  \node[store,left=10mm of tc,text width=2.6cm] (bin) {\file{data/train.bin}\\ \file{data/val.bin}\\ \scriptsize flat uint16 stream};

  \node[box,below=12mm of bin,text width=3.0cm] (train) {\code{train.py}\\ \scriptsize AdamW, cosine LR,\\ bf16, DDP, resume};
  \node[box,left=10mm of train,text width=2.6cm] (model) {\code{model/}\\ \scriptsize RMSNorm, RoPE,\\ SwiGLU, tied emb};
  \node[store,right=10mm of train,text width=3.1cm] (runs) {\file{runs/<preset>/}\\ \scriptsize ckpt.pt, model.pt,\\ metrics.jsonl};
  \node[box,right=10mm of runs,text width=3.0cm] (ev) {\code{eval/} + \code{sample.py}\\ \scriptsize perplexity, script mix,\\ cloze, tokenizer cmp};

  \draw[flow] (wiki.east) -- (src.west);
  \draw[flow] (leip.east) -- (src.west);
  \draw[flow] (cc.east) -- (src.west);
  \draw[dflow] (wd.east) -- (src.west);
  \draw[flow] (src) -- (norm);
  \draw[flow] (norm) -- (dd);
  \draw[flow] (dd) -- (sp);
  \draw[flow] (sp) -- (jsonl);
  \draw[flow] (jsonl) -- (tt) node[lbl,midway,right] {train once};
  \draw[flow] (tt) -- (tokj);
  \draw[flow] (tokj) -- (tc);
  \draw[flow] (jsonl.west) to[out=180,in=90] (tc.north);
  \draw[flow] (tc) -- (bin);
  \draw[flow] (bin) -- (train) node[lbl,midway,right] {memory-mapped};
  \draw[flow] (model) -- (train);
  \draw[flow] (train) -- (runs);
  \draw[flow] (runs) -- (ev);
\end{tikzpicture}}
\caption{End-to-end architecture. Solid arrows are paths that have been executed and
produced the committed artifacts. The Wikidata source is wired and unit-tested but has
not been used in any training run.}
\label{fig:arch}
\end{figure}

\subsection{The repository, file by file}

\begin{figure}[H]
\centering
\begin{forest} dirtree
[urdu-slm
  [pipeline/ {\rmfamily\itshape raw text $\rightarrow$ token stream}
    [sources.py {\rmfamily\itshape readers: wiki dump, Leipzig tar, plain text, Wikidata}]
    [normalize.py {\rmfamily\itshape NFC, variant folding, script-ratio language ID}]
    [dedup.py {\rmfamily\itshape exact blake2b + vectorized MinHash LSH}]
    [common.py {\rmfamily\itshape ShardWriter, Stage resume markers, iter\_jsonl}]
    [build\_corpus.py {\rmfamily\itshape orchestrates extract $\rightarrow$ dedup $\rightarrow$ split}]
    [train\_tokenizer.py {\rmfamily\itshape 32k byte-level BPE}]
    [tokenize\_corpus.py {\rmfamily\itshape encode to uint16 .bin, upsample, filter}]
    [fetch\_wikidata.py {\rmfamily\itshape SPARQL fact triples (standalone)}]
  ]
  [model/
    [config.py {\rmfamily\itshape ModelConfig dataclass + 4 size presets}]
    [transformer.py {\rmfamily\itshape the entire network, 186 lines}]
  ]
  [train.py {\rmfamily\itshape the pretraining loop}]
  [sample.py {\rmfamily\itshape generate continuations from a checkpoint}]
  [eval/
    [\_shared.py {\rmfamily\itshape load\_model helper}]
    [perplexity.py {\rmfamily\itshape held-out loss on a .bin split}]
    [script\_mix.py {\rmfamily\itshape share of generated chars in Arabic script}]
    [cloze.py {\rmfamily\itshape few-shot fill-in-the-blank harness}]
    [tokenizer\_compare.py {\rmfamily\itshape ours vs GPT-2 compression}]
    [plot\_loss.py {\rmfamily\itshape metrics.jsonl $\rightarrow$ PNG curve}]
  ]
  [configs/ {\rmfamily\itshape tiny.yaml small.yaml base.yaml medium.yaml}]
  [evaldata/cloze.jsonl {\rmfamily\itshape 95 hand-written Urdu items}]
  [tests/ {\rmfamily\itshape test\_model.py (5), test\_pipeline.py (8)}]
  [runs/ {\rmfamily\itshape tiny/, base/, base\_v2/ + metrics.jsonl (LFS weights)}]
  [data/ {\rmfamily\itshape val.bin, corpus\_stats.json, sample/ (train.bin gitignored)}]
  [docs/ {\rmfamily\itshape GPU\_RUN.md, DATA\_LICENSES.md, explainers/}]
]
\end{forest}
\caption{Directory structure with each file's job. The whole project is about 2,800
lines of Python including tests and docs.}
\label{fig:tree}
\end{figure}

\begin{keyidea}{The organizing principle}
Each stage writes its output to disk and never holds the corpus in memory. Sources
\emph{stream} (Python generators, one document at a time), writers \emph{shard} (fixed
size JSONL files), stages \emph{mark themselves done} (a marker file), and training
\emph{memory-maps} the token stream. A 1.7GB corpus is processed on a machine that never
allocates 1.7GB for it. This is the single most consistent design decision in the
repository.
\end{keyidea}

\newpage
% =====================================================================
\section{Part III: The data pipeline}
\label{sec:pipeline}

The \file{pipeline/} package answers one question: how do you turn a 455MB compressed
Wikipedia dump and a pile of news sentences into a clean, deduplicated, tokenized stream
of numbers, on a laptop, without ever loading it all into memory?

\subsection{\code{sources.py}: one reader per format}

Every source is a Python \emph{generator} yielding one dictionary per document, shaped
\code{\{text, source\}}. This uniformity is what lets the rest of the pipeline stay
source-agnostic.

\begin{jargon}{Generator (Python)}
A function that produces values one at a time with \code{yield} instead of building a
whole list. The caller consumes them in a loop, and only one value exists in memory at
a time. This is the mechanism behind the streaming claim: \code{read\_wiki\_dump} on a
455MB dump never holds more than one article.
\end{jargon}

\subsubsection{\code{read\_wiki\_dump}: parsing a MediaWiki dump}

A Wikipedia dump is a single enormous XML file, bz2-compressed. The reader uses
\code{ET.iterparse}, which fires an event as each element \emph{closes}, rather than
building a tree of the whole document.

\begin{lstlisting}[language=Python,caption={The streaming discipline in read\_wiki\_dump}]
for event, elem in ET.iterparse(fh, events=("end",)):
    tag = elem.tag.split("}")[-1]
    if tag == "title":   title = elem.text
    elif tag == "ns":    ns = elem.text
    elif tag == "text":  text = elem.text
    elif tag == "page":
        if ns == "0" and text and title and not _REDIRECT_RE.match(text):
            body = _strip_wikitext(text)
            ...
        title, text, ns = None, None, None
        elem.clear()          # <- releases the parsed element's memory
\end{lstlisting}

Three details carry real weight:

\begin{itemize}
\item \code{elem.clear()} is what makes this constant-memory. Without it,
      \code{iterparse} still accumulates every parsed element and the process grows to
      the size of the dump.
\item \code{tag.split("\}")[-1]} strips the XML namespace. MediaWiki tags do not arrive
      as a bare \code{title}; they arrive with the schema URL wrapped in braces in front
      of the tag name, as shown below. Splitting on the closing brace and taking the last
      piece discards the namespace and leaves the plain tag name.
\begin{lstlisting}[caption={The namespaced tag that the split has to strip}]
{http://www.mediawiki.org/xml/export-0.10/}title   ->   title
\end{lstlisting}
\item \code{ns == "0"} keeps only article-namespace pages, dropping Talk pages, User
      pages, Templates and Categories. \code{\_REDIRECT\_RE} then drops redirect stubs,
      matching \code{\#REDIRECT} in English \emph{and} two Urdu spellings of the same
      directive.
\end{itemize}

\code{\_strip\_wikitext} hands the wikitext to \code{mwparserfromhell}, an external
library that understands MediaWiki markup, and calls \code{strip\_code}. It then drops
any surviving line starting with a pipe, exclamation mark, or brace, which are table and
infobox leftovers the parser does not fully clean.

\begin{jargon}{Wikitext}
Wikipedia's own markup language: \code{[[links]]}, \code{\{\{templates\}\}}, table
syntax. It is famously irregular, which is why parsing it is delegated to a dedicated
library rather than attempted with regular expressions.
\end{jargon}

\subsubsection{The other three readers}

\begin{itemize}
\item \code{read\_leipzig\_tar} opens a \file{.tar.gz}, scans members for the one ending
      \file{-sentences.txt}, and splits each line on the first tab, since Leipzig's
      format is \code{<id>\textbackslash t<sentence>}. If no such member exists it
      returns silently, yielding nothing.
\item \code{read\_plain\_text} handles the CC-100 layout: consecutive non-empty lines
      form one document, a blank line ends it. It transparently opens \file{.xz},
      \file{.gz} or plain files, and passes \code{errors="ignore"} so a single malformed
      byte in a 9-million-document crawl cannot kill the run.
\item \code{read\_wikidata\_facts} reads the JSONL written by
      \code{pipeline/fetch\_wikidata.py} and renders each subject/relation/object triple
      into a declarative Urdu sentence using two fixed templates per relation.
\end{itemize}

\begin{keyidea}{Why the Wikidata templates alternate by index, not at random}
\code{template = templates[i \% len(templates)]} picks the phrasing from the row's line
number. Using \code{random.choice} would give the same lexical variety but a different
corpus on every run, which would silently break reproducibility. Deterministic
alternation buys the variety without the cost. This is a small decision, well made.
\end{keyidea}

\subsection{\code{normalize.py}: making Urdu canonical}
\label{sec:normalize}

This module is where the project's Urdu-specific knowledge is concentrated, and it is
the part a generic "train a GPT" tutorial does not have at all.

\subsubsection{The folding map}

Section~\ref{sec:whyurdu} explained why encoding variants must be collapsed. The code
does it with a compiled alternation regex over a nine-entry map. Because the actual
characters cannot be printed here, Table~\ref{tab:fold} gives the map by codepoint and
name.

\begin{table}[H]
\centering
\small
\begin{tabular}{llll}
\toprule
\textbf{From} & \textbf{Name} & \textbf{To} & \textbf{Why} \\
\midrule
U+0643 & Arabic Kaf            & U+06A9 (Urdu Keheh) & same letter, Arabic spelling \\
U+064A & Arabic Yeh            & U+06CC (Farsi Yeh)  & same letter, Arabic spelling \\
U+0649 & Alef Maksura          & U+06CC (Farsi Yeh)  & same letter, Arabic spelling \\
U+06C0 & Heh with Hamza Above  & U+06C1 (Heh Goal)   & common variant \\
U+200B & Zero-width space      & (removed)           & invisible, no meaning \\
U+200D & Zero-width joiner     & (removed)           & invisible, no meaning \\
U+FEFF & BOM                   & (removed)           & invisible, no meaning \\
U+00A0 & Non-breaking space    & U+0020 (space)      & normalize whitespace \\
U+0640 & Tatweel / kashida     & (removed)           & decorative elongation only \\
\bottomrule
\end{tabular}
\caption{The variant-folding map from \code{\_FOLD}, given by codepoint because the
glyphs cannot be typeset here.}
\label{tab:fold}
\end{table}

\begin{keyidea}{The zero-width non-joiner is deliberately kept}
U+200C, the zero-width non-joiner, is \emph{not} in the removal list, while its
neighbours U+200B and U+200D are. This is not an oversight. In Urdu the ZWNJ carries
real orthographic meaning: it stops two letters from joining where joining would change
the word. Deleting it would corrupt text. \code{tests/test\_pipeline.py} pins this
exact behaviour with an assertion that a ZWNJ survives normalization while a tatweel
next to it does not. That is a well-chosen test.
\end{keyidea}

The module's docstring also records a decision worth defending out loud: \emph{harakat}
(short-vowel diacritics) are \textbf{not} stripped. They are rare in running Urdu but
meaningful when present, for instance in religious quotation and dictionaries, so
removing them would corrupt that minority of text. Many pipelines strip them by reflex.

\subsubsection{NFC normalization}

\begin{jargon}{Unicode normalization, NFC}
The same visible character can sometimes be encoded either as one codepoint or as a base
character plus a combining mark. NFC (Normalization Form Composed) rewrites text so the
composed form always wins. Without it, two identical-looking strings can compare as
different, which would defeat exact deduplication and split the tokenizer's vocabulary.
\code{normalize()} applies NFC \emph{first}, then folds variants.
\end{jargon}

\subsubsection{Language identification without a model}

\code{arabic\_ratio(text)} counts what share of the \emph{letters} (characters where
\code{c.isalpha()} is true, so digits and punctuation are excluded) fall in one of five
Unicode ranges: Arabic (U+0600 to U+06FF), Arabic Supplement, Arabic Extended-A, and the
two Presentation Forms blocks.

\begin{keyidea}{The reasoning behind skipping fastText}
The obvious tool is fastText's \code{lid.176} language-ID model. It is 126MB and adds a
download plus a dependency. The module argues that Urdu shares its script only with other
Perso-Arabic languages, none of which appear at volume in these sources (a Wikipedia
dump already filtered to \code{ur}, and Leipzig corpora already tagged \code{urd}), so
the real contaminant is embedded English, which a script ratio catches directly. For
\emph{these} sources that reasoning holds. The README is candid that a mixed-script web
crawl would need real token-level language ID, and CC-100 (added in v2) is exactly such
a crawl.
\end{keyidea}

\code{keep\_line()} then applies four filters in order, all cheap tests first:

\begin{table}[H]
\centering
\small
\begin{tabular}{lll}
\toprule
\textbf{Filter} & \textbf{Threshold} & \textbf{Catches} \\
\midrule
\code{len(text) < MIN\_CHARS}                & 40   & fragments, headings \\
\code{len(text.split()) < MIN\_WORDS}        & 6    & captions, list items \\
\code{arabic\_ratio(text) < MIN\_ARABIC\_RATIO} & 0.65 & English and other-script lines \\
alpha share of length $< 0.5$                & 0.5  & tables, infobox leftovers, digit soup \\
\bottomrule
\end{tabular}
\caption{The boilerplate and language filter, \code{normalize.keep\_line}.}
\end{table}

The README records that writing the test for this surfaced a bug in the author's own
expectation, not in the code: a mixed Urdu/English phrase is 4 Arabic letters of 16, so
0.25 rather than the 0.5 first assumed. The test now pins the real value with a range
assertion. That is the test doing its job.

\subsection{\code{dedup.py}: the most interesting code in the repository}
\label{sec:dedup}

Deduplication matters here more than in a typical corpus, for a reason specific to this
language: Urdu Wikipedia is largely bot-generated stub articles. The pipeline extracted
2.46 million paragraphs from it, far more than the article count would suggest, and a
large share are near-identical templated text. Training on them would teach the model to
recite templates.

\subsubsection{Exact duplicates}

\code{exact\_key} hashes the normalized text with blake2b truncated to 8 bytes, and a
Python \code{set} holds the keys. Storing a 64-bit digest instead of the text is what
keeps 3.7 million documents' worth of keys in memory.

\begin{jargon}{Hash function}
A function that maps data of any size to a fixed-size number, such that the same input
always gives the same output and different inputs almost never collide. Comparing hashes
is a cheap stand-in for comparing the full text.
\end{jargon}

\subsubsection{Near duplicates: MinHash and LSH}

Exact hashing misses the real problem: two paragraphs differing by one character have
completely unrelated hashes. That is what MinHash is for.

\begin{jargon}{Shingle}
A short overlapping substring. \code{\_shingles(text, k=5)} strips all whitespace and
takes every 5-character window, deduplicated into a set. "abcdefg" gives
\code{\{abcde, bcdef, cdefg\}}. Two documents that share most of their shingles are
near-duplicates.
\end{jargon}

\begin{jargon}{Jaccard similarity}
The size of the intersection of two sets divided by the size of their union. 1.0 means
identical sets, 0.0 means no overlap. Here the sets are the two documents' shingles, and
the pipeline's threshold is 0.8.
\end{jargon}

\begin{jargon}{MinHash}
A trick for estimating Jaccard similarity without comparing the sets. Apply a random
permutation to the universe of possible shingles and record only the \emph{minimum}
element of each set under it. The probability that two sets agree on that minimum is
exactly their Jaccard similarity. Repeat with 64 different permutations and you get a
64-number \emph{signature}; the share of positions where two signatures agree estimates
their Jaccard. Constant size, regardless of document length.
\end{jargon}

\begin{jargon}{LSH (Locality-Sensitive Hashing)}
Even with signatures, comparing every new document against every kept one is quadratic.
LSH splits each 64-number signature into bands, hashes each band into a bucket, and only
compares documents that collide in at least one band. Similar documents collide with high
probability; dissimilar ones almost never do. This turns "compare against everything"
into "look in a few buckets".
\end{jargon}

\begin{figure}[H]
\centering
\resizebox{\textwidth}{!}{
\begin{tikzpicture}[node distance=8mm]
  \node[box,text width=2.6cm] (doc) {one document\\ \scriptsize normalized paragraph};
  \node[box,right=of doc,text width=2.7cm] (sh) {\code{\_shingles(k=5)}\\ \scriptsize set of 5-char windows,\\ whitespace stripped};
  \node[box,right=of sh,text width=2.9cm] (xx) {\code{xxhash.xxh32}\\ \scriptsize every shingle to a\\ 32-bit int, vectorized\\ into one numpy array};
  \node[box,right=of xx,text width=3.5cm] (perm) {affine permutations\\ \code{(outer(hv,A)+B) \% prime}\\ \scriptsize one (E x 64) matrix};
  \node[box,below=10mm of perm,text width=3.2cm] (min) {\code{.min(axis=0)}\\ \scriptsize 64-number signature};
  \node[dec,left=12mm of min,text width=2.0cm] (q) {\code{lsh.query(m)}\\ hit?};
  \node[box,left=14mm of q,text width=2.5cm,draw=warnred] (drop) {drop\\ \scriptsize \code{n\_near += 1}};
  \node[box,below=9mm of q,text width=2.5cm,draw=okgreen] (keep) {keep + insert\\ \scriptsize into the LSH index};

  \draw[flow] (doc) -- (sh);
  \draw[flow] (sh) -- (xx);
  \draw[flow] (xx) -- (perm);
  \draw[flow] (perm) -- (min);
  \draw[flow] (min) -- (q);
  \draw[flow] (q) -- node[lbl,above] {yes} (drop);
  \draw[flow] (q) -- node[lbl,right] {no} (keep);
  \node[extbox,text width=6.6cm,below=9mm of keep]
    {\scriptsize The two numpy operations above (the affine permutations and the min)\\
     replace datasketch's per-shingle Python loop.\\
     Measured on real Wikipedia paragraphs: $\sim$250 docs/s to $\sim$2400 docs/s.};
\end{tikzpicture}}
\caption{The near-duplicate path for a single document. Exact blake2b dedup runs first
and short-circuits before any of this.}
\label{fig:dedup}
\end{figure}

\subsubsection{The vectorization, and why it was necessary}

This is the repository's best engineering story, and it is worth being able to tell it
precisely.

\begin{lstlisting}[language=Python,caption={The hot path, rewritten}]
_seed = MinHash(num_perm=64, seed=1, scheme="affine32")
_A, _B = _seed.permutations          # affine coefficients, shape (64,)

def build_minhash(text, num_perm=_NUM_PERM):
    m = MinHash(num_perm=num_perm, seed=1, permutations=(_A, _B), scheme="affine32")
    shs = _shingles(text)
    if not shs:
        return m
    hv = np.fromiter((xxhash.xxh32(s).intdigest() for s in shs),
                     dtype=np.uint64, count=len(shs))
    phv = (np.outer(hv, _A) + _B) % _MERSENNE      # (E, 64) in one op
    phv = np.bitwise_and(phv, _MAX_HASH)
    m.hashvalues = phv.min(axis=0).astype(np.uint64)   # (64,)
    return m
\end{lstlisting}

The library's own API is \code{MinHash.update(shingle)}, called once per shingle. That
loops in Python and runs a tiny numpy operation per call, so per-call overhead dominates
and the work never reaches numpy's fast path. Benchmarked on real Wikipedia paragraphs it
managed about 250 documents per second, which extrapolates to over two hours for the full
corpus.

The rewrite computes the entire signature in two array operations. \code{np.outer(hv,
\_A)} builds the whole (shingles $\times$ 64) matrix of permuted values at once, the
modulo and mask apply elementwise, and \code{.min(axis=0)} collapses it to the 64-number
signature. Measured throughput went to about 2400 documents per second, roughly a 10x
speedup.

\begin{keyidea}{Why swapping the hash function is legal}
This is the insight that makes the rewrite safe. MinHash signatures do not need to match
any external standard; they only need to be \emph{internally consistent}. Every document
must be hashed the same way as every other, and that is the whole requirement. That
freedom permits replacing datasketch's sha1 element hash with the much faster xxhash and
writing \code{hashvalues} directly, as long as all documents share one scheme and one set
of permutations. Reusing the seed MinHash's own \code{\_A, \_B} guarantees exactly that,
which is why the LSH banding on top continues to work unchanged.
\end{keyidea}

xxhash is also chosen over Python's builtin \code{hash()} for a specific reason:
\code{hash()} on strings is salted per process, so it returns different values on
different runs and would make the corpus irreproducible.

The README also records a real integration bug: datasketch 2.0 changed the MinHash
constructor so that passing \code{permutations=} without a \code{scheme} raises
\code{ValueError: scheme must be specified explicitly}. The fix was to read the seed
MinHash's scheme (\code{affine32}) and pass it explicitly everywhere.

\begin{honest}{A latent trap in \code{build\_minhash}'s signature}
\code{build\_minhash(text, num\_perm=\_NUM\_PERM)} takes \code{num\_perm} as a
parameter, but \code{\_A} and \code{\_B} are module-level constants built once at
\code{num\_perm=64}. Passing anything other than 64 would construct a MinHash claiming a
different permutation count while handing it 64 coefficients. Nothing in the repository
does this (\code{Deduper} defaults to 64 and passes its own \code{self.num\_perm}, also
64), so it is not a live bug. But the parameter advertises a flexibility the function
does not have. Either wire it through properly or drop the argument.
\end{honest}

\subsubsection{\code{Deduper}: the two filters composed}

\code{Deduper.keep(text)} runs exact first (cheap, short-circuits), then near. It keeps
three counters, \code{n\_exact}, \code{n\_near}, \code{n\_kept}, and \code{stats()}
derives the dedup rate from them. Documents are inserted into the LSH under a
monotonically increasing integer id, which is never used to look anything up: only
\code{query()}'s truthiness matters, not which document was hit.

\begin{honest}{The LSH index lives entirely in RAM}
\code{MinHashLSH(threshold=0.8, num\_perm=64)} with datasketch's default storage is an
in-memory index, and \code{seen\_exact} is an in-memory set. For 2.5 million (v1) and
13.3 million (v2) kept documents, both grow with the corpus. This is the real memory
ceiling of the whole pipeline, and it is the one place where the otherwise-strict
streaming discipline does not hold. The README names this itself under "What I'd do
differently" and proposes persisting signatures to disk so the stage becomes
independently resumable and the Jaccard threshold could be retuned without re-hashing.
The fix is not implemented; the honest position is that the v2 run's memory footprint
worked on the machine it was run on.
\end{honest}

\subsection{\code{common.py}: the plumbing that makes stages resumable}

Three small pieces, each doing exactly one job.

\begin{itemize}
\item \code{iter\_jsonl(path)} yields one parsed object per non-blank line. Streaming,
      naturally.
\item \code{ShardWriter} writes records to \file{prefix-00000.jsonl},
      \file{prefix-00001.jsonl}, rolling to a new file every \code{docs\_per\_shard}
      records. It opens lazily (no empty file if nothing is ever written) and supports
      the \code{with} statement. \code{ensure\_ascii=False} keeps Urdu readable in the
      shard rather than escaped to \code{\textbackslash u06cc} noise.
\item \code{Stage} implements resume with the simplest mechanism that works: a marker
      file named \file{.<stage>.done} that holds the stage's own stats as JSON.
      \code{is\_done()} is an existence check.
\end{itemize}

\begin{jargon}{JSONL}
"JSON Lines": a text file with one complete JSON object per line. Unlike a single JSON
array, it can be read one record at a time and appended to without rewriting, which is
why every streaming pipeline uses it.
\end{jargon}

\begin{honest}{The resume marker is coarser than it looks}
The marker is written only after a stage \emph{completes}. If extraction dies 90\%
through a nine-hour CC-100 pass, no marker exists, and the next run redoes all of it.
"Resumable" here means "skip stages that finished", not "continue from where it broke".
That is still valuable (the three stages are hours apart) but it is a weaker guarantee
than the word usually implies, and worth stating plainly rather than letting an
interviewer discover it.
\end{honest}

\subsection{\code{build\_corpus.py}: the orchestrator}

\code{main()} wires three stage functions in order and writes
\file{data/corpus\_stats.json} at the end. The design is worth noting: every stage is a
plain function taking \code{(args, work)}, checks its own marker, and returns its stats
dict. There is no framework.

\begin{itemize}
\item \textbf{extract}: for each configured source, stream documents, \code{normalize()}
      the text, split into paragraphs on newlines, and write those that pass
      \code{keep\_line()} to a per-source shard. Note the granularity change here: a
      source yields \emph{documents} (whole articles) but the pipeline stores
      \emph{paragraphs}. That is why 1.73M Wikipedia articles produce 2.46M kept
      paragraphs.
\item \textbf{dedup}: read every extract shard through one \code{Deduper}, writing
      survivors to \file{clean/}. Deduplication is corpus-wide, not per-source, so a
      paragraph appearing in both Wikipedia and CC-100 is caught.
\item \textbf{split}: read the clean shards and send each record to val when
      \code{rng.random()} falls below \code{args.val\_frac} (default 0.01), and to train
      otherwise. The generator is seeded at 1337, so the split is reproducible.
\end{itemize}

\begin{honest}{The split is a per-document coin flip, not a guaranteed 1\%}
\code{stage\_split} draws an independent random number per document. The expected val
share is exactly 1\%, and at 2.5 million documents the variance is negligible (the
observed v1 split, 25,106 val of 2,515,082, is 0.998\%). It is deterministic given the
seed, which is what matters for reproducibility. But it is a Bernoulli draw, not a
partition, so the count is not exactly 1\% and "deterministic 1\% holdout" is a
description of the mechanism's expectation rather than its guarantee. This is a nitpick,
not a defect.
\end{honest}

\subsection{What the pipeline actually produced}

Figure~\ref{fig:funnel} traces the v1 corpus, whose real numbers are the committed
\file{data/corpus\_stats.json}. Every figure below was read out of that file, not out of
the README.

\begin{figure}[H]
\centering
\resizebox{0.94\textwidth}{!}{
\begin{tikzpicture}[node distance=6mm]
  \node[box,text width=11cm,minimum height=11mm] (a) {\textbf{3,028,010 raw units in}\\ \scriptsize 1,728,010 Wikipedia articles + 1,000,000 + 300,000 Leipzig sentences};
  \node[box,text width=9.4cm,below=of a,minimum height=11mm] (b) {\textbf{3,695,487 paragraphs kept after filtering}\\ \scriptsize wiki 2,459,293 (articles split into paragraphs) + 960,189 + 276,005};
  \node[box,text width=7.4cm,below=of b,minimum height=11mm] (c) {\textbf{2,515,082 documents after dedup}\\ \scriptsize 871,272 exact + 309,133 near removed = 31.94\%};
  \node[box,text width=5.6cm,below=of c,minimum height=11mm] (d) {\textbf{438.5M characters kept}\\ \scriptsize 2,489,976 train / 25,106 val};
  \node[box,text width=4.0cm,below=of d,minimum height=11mm,draw=okgreen] (e) {\textbf{112.8M tokens}\\ \scriptsize 111.7M train + 1.1M val};
  \draw[flow] (a) -- (b);
  \draw[flow] (b) -- (c);
  \draw[flow] (c) -- (d);
  \draw[flow] (d) -- (e);
  \node[lbl,anchor=west,text width=5.2cm] at ($(b.east)+(0.3,0)$) {\scriptsize expands: one article\\ becomes many paragraphs};
  \node[lbl,anchor=west,text width=5.2cm] at ($(c.east)+(1.3,0)$) {\scriptsize the bot-stub collapse\\ this stage exists for};
  \node[lbl,anchor=west,text width=5.2cm] at ($(e.east)+(2.1,0)$) {\scriptsize 4.078 chars/token,\\ measured};
\end{tikzpicture}}
\caption{The v1 corpus funnel. All figures are from the committed
\file{data/corpus\_stats.json}, verified by reading it during the writing of this
document.}
\label{fig:funnel}
\end{figure}

\begin{honest}{The committed corpus stats describe v1, but the shipped model is v2}
\file{data/corpus\_stats.json} contains exactly three sources (wikipedia, leipzig 2016,
leipzig 2011) and the v1 numbers above. There is \textbf{no} committed stats file for
the v2 corpus, the one that CC-100 was added to and that \file{runs/base\_v2} (the
project's headline model) was actually trained on. The v2 table in the README (11.0M
CC-100 documents kept, 37.5\% dedup rate, 13.3M train docs) is therefore the \emph{only}
record of that corpus, and unlike the v1 numbers it cannot be checked against a committed
artifact. It is not doubted here; the point is that it rests on the README's word where
v1's rests on a file. Re-running the pipeline would regenerate the file and overwrite the
v1 stats, which may be exactly why it was not re-run.
\end{honest}

\newpage
% =====================================================================
\section{Part IV: The tokenizer}
\label{sec:tokenizer}

This is the part of the project that most directly answers "why build this at all
instead of fine-tuning something". Section~\ref{sec:whyurdu} argued from first
principles that an English tokenizer must handle Urdu badly. This part measures it.

\subsection{\code{train\_tokenizer.py}: 56 lines, several real decisions}

\begin{lstlisting}[language=Python,caption={The whole configuration of the tokenizer}]
tok = Tokenizer(models.BPE(unk_token="<|unk|>"))
tok.pre_tokenizer = pre_tokenizers.ByteLevel(add_prefix_space=False)
tok.decoder = decoders.ByteLevel()
trainer = trainers.BpeTrainer(
    vocab_size=32000,
    min_frequency=2,
    special_tokens=["<|endoftext|>", "<|pad|>", "<|unk|>"],
    initial_alphabet=pre_tokenizers.ByteLevel.alphabet(),
    show_progress=True,
)
tok.train_from_iterator(corpus_lines([data / "train"]), trainer=trainer)
tok.post_processor = processors.ByteLevel(trim_offsets=False)
\end{lstlisting}

Unpacking each choice:

\begin{jargon}{Byte-level BPE}
BPE whose base alphabet is the 256 possible byte values, not the set of characters seen
in training. Since every possible input encodes to bytes, and every byte is in the
vocabulary, \textbf{no input can ever be out-of-vocabulary}. A character-level BPE would
have to emit an unknown token for a character it never saw; a byte-level one always has
a fallback, at worst spending one token per byte.
\end{jargon}

\begin{itemize}
\item \code{initial\_alphabet=ByteLevel.alphabet()} is what guarantees that property. It
      seeds the vocabulary with all 256 bytes \emph{before} training, so even bytes that
      never appear in Urdu text (an emoji, a Cyrillic letter) still have an id. Without
      it, only bytes observed in the corpus would be present.
\item \code{min\_frequency=2} refuses to spend a vocabulary slot on a merge seen once.
\item \code{add\_prefix\_space=False} is a deliberate departure from GPT-2's default.
      GPT-2 prepends a space to the input so that "word" and " word" tokenize alike.
      That convention is built around English's space-delimited orthography; this
      tokenizer declines it.
\item \textbf{The tokenizer is trained on \code{data/train} only.} \code{shard\_dirs =
      [data / "train"]}, with the val directory pointedly excluded. This matters more
      than it looks: a tokenizer trained on the validation set would have merges shaped
      by text the model is later evaluated on, a subtle leak. Getting this right is a
      mark of care.
\end{itemize}

\begin{jargon}{Special tokens}
Ids reserved for control rather than text. \code{<|endoftext|>} marks a document
boundary and is appended after every document during tokenization.  \code{<|pad|>} would
fill short sequences to equal length, and \code{<|unk|>} would represent an unknown
symbol.
\end{jargon}

\begin{honest}{Two of the three special tokens are never used}
\code{<|pad|>} is declared and given a vocabulary slot, but nothing in the repository
ever pads: \code{TokenLoader} slices fixed-length blocks out of one continuous token
stream, so every sequence is exactly \code{block\_size} long by construction and padding
is structurally unnecessary. \code{<|unk|>} is passed to \code{models.BPE(unk\_token=...)}
but a byte-level BPE with a full 256-byte initial alphabet can never emit it, by the
very design described above. Only \code{<|endoftext|>} does real work. This costs two
vocabulary slots out of 32,000, which is nothing, and both are conventional insurance.
It is worth knowing they are inert rather than being asked in an interview what
\code{<|pad|>} is for and discovering the answer is "nothing".
\end{honest}

\subsection{\code{tokenize\_corpus.py}: text to a flat array of numbers}

Once the tokenizer is frozen, the corpus is encoded once, ahead of time, into a single
flat binary file per split. Nothing is tokenized during training.

\begin{lstlisting}[language=Python,caption={The core of encode\_split}]
for rec in iter_jsonl(shard):
    src = rec.get("source")
    if only and src not in only:        # --filter-source
        continue
    ids = tok.encode(rec["text"]).ids
    ids.append(eot_id)                  # every doc ends with <|endoftext|>
    reps = upsample.get(src, 1) if upsample else 1
    buf.extend(ids * reps)              # --upsample wikipedia:3
arr = np.array(buf, dtype=np.uint16)
arr.tofile(fh)
\end{lstlisting}

\begin{jargon}{uint16}
An unsigned 16-bit integer: whole numbers from 0 to 65,535. The vocabulary is 32,000, so
every token id fits, and each costs 2 bytes instead of the 4 an int32 would take. On the
862M-token v2 corpus that is the difference between a 1.7GB file and a 3.4GB one. The
choice is only safe because the vocabulary is under 65,536, which is a constraint worth
knowing about before someone raises \code{vocab\_size}.
\end{jargon}

\begin{jargon}{Memory-mapping}
Handing the operating system a file and asking it to make the bytes appear as if they
were an array in memory, without reading them in. Pages load on demand and the OS
evicts them under pressure. \code{np.memmap(bin\_path, dtype=np.uint16, mode="r")} is
how a 1.7GB corpus is trained on without a 1.7GB allocation.
\end{jargon}

\subsubsection{Upsampling and filtering, and why both exist}

The v2 corpus is dominated by CC-100, a web crawl, by sheer volume. Wikipedia is smaller
but cleaner. Left alone, the model would spend most of its training reading crawl text.
Two flags address this:

\begin{itemize}
\item \code{--upsample wikipedia:3} repeats every Wikipedia document three times in
      \file{train.bin}, so a smaller high-quality source is not swamped. This produced
      the 862.4M-token main mix.
\item \code{--filter-source wikipedia} writes a Wikipedia-only stream, used for a short
      \emph{annealing} phase at the end of training. This produced the 70.8M-token
      \file{data\_anneal/train.bin}.
\end{itemize}

\begin{keyidea}{Annealing, in this project's sense}
Training first on the big mixed corpus, then finishing on a small high-quality slice.
The idea is that the final gradient steps have outsized influence on the model's habits,
so ending on clean encyclopedic prose biases the finished model toward it. The v2 run did
exactly this, and Section~\ref{sec:runs} shows the switch is plainly visible in the loss
log.
\end{keyidea}

\begin{keyidea}{Validation is never upsampled or filtered}
\code{upsample=upsample if train else None} and \code{only=only if train else None}. The
val split is always the full, unweighted distribution. If val were upsampled to match
train, validation loss would measure performance on a distorted distribution and would
not be comparable across runs with different mixes. This is quietly correct, and it is
also why \file{data\_anneal/val.bin} and \file{data/val.bin} both hold the same 7.3M
tokens.
\end{keyidea}

\begin{honest}{Upsampling repeats a document's tokens back-to-back inside one context}
\code{buf.extend(ids * reps)} is Python list repetition, which lays the three copies
down \textbf{contiguously}: \code{[doc, doc, doc]} in the stream, not three copies
scattered across the corpus. \code{TokenLoader} then slices blocks of 1024 tokens at
fixed offsets from that stream. A typical kept Wikipedia paragraph is a few hundred
characters, so at the measured 4.078 chars/token it is well under 100 tokens, and all
three copies of it comfortably fit inside a single 1024-token training block.

The consequence is that the model repeatedly sees a context in which the literal correct
continuation is a verbatim copy of text 70 tokens earlier. That is a strong, easily
learned signal, and it rewards copying rather than modeling. This is a different thing
from what upsampling is meant to do, which is to make Wikipedia's \emph{distribution}
count for more across the corpus.

The fix is cheap: repeat at the shard or document-list level, or interleave the copies,
rather than repeating the id list in place. Whether it measurably hurt the v2 model is
not established here and cannot be established without a retrain. It is stated as a real
design defect that survived into the run that produced the headline model, and it is
probably the single most substantive code finding in this document.
\end{honest}

\subsection{The 5.23x claim, verified}
\label{sec:tokclaim}

The README's headline claim is that GPT-2's tokenizer spends 5.23x more tokens on the
same Urdu text. This was checked by running the code rather than trusting the table.

\code{eval/tokenizer\_compare.py} builds a sample by streaming records out of
\file{data/val} until it has \code{--max-chars} characters, encodes it with both
tokenizers, and reports chars-per-token for each. The comparison is deterministic: no
sampling, no seed, so it must reproduce exactly.

\begin{keyidea}{Verified on this machine, 2026-07-16}
A real GPT-2 \file{tokenizer.json} was found in the local Hugging Face cache and the
comparison was run against it. Every number in the README's table reproduced
\textbf{exactly}, to the digit:
\end{keyidea}

\begin{lstlisting}[caption={Real output, reproducing the README's table exactly}]
$ python -m eval.tokenizer_compare --gpt2 <cached gpt2 tokenizer.json> \
      --max-chars 800000
sample chars: 806495
ours (32k Urdu BPE): 197758 tokens, 4.078 chars/token
gpt2: 1034831 tokens, 0.779 chars/token
compression gain: 5.23x fewer tokens with our tokenizer
\end{lstlisting}

\begin{figure}[H]
\centering
\begin{tikzpicture}
\begin{axis}[
  width=13cm, height=6.0cm,
  xbar, y=9mm, bar width=6mm,
  xmin=0, xmax=1380000,
  xlabel={tokens spent on the same 806,495 characters of Urdu (lower is better)},
  symbolic y coords={gpt2,ours},
  ytick=data,
  yticklabels={GPT-2 (50k English BPE), ours (32k Urdu BPE)},
  nodes near coords, nodes near coords align={horizontal},
  nodes near coords style={font=\footnotesize, text=inkblue,
    /pgf/number format/fixed, /pgf/number format/precision=0,
    /pgf/number format/1000 sep={,}},
  every axis plot/.append style={fill=accent!55, draw=inkblue},
  tick label style={font=\small}, label style={font=\small},
  xmajorgrids, grid style={linegrey},
  scaled x ticks=false, xtick={0,250000,500000,750000,1000000},
  xticklabel style={font=\small, /pgf/number format/fixed,
    /pgf/number format/1000 sep={,}},
  enlarge y limits=0.55,
]
\addplot coordinates {(1034831,gpt2) (197758,ours)};
\end{axis}
\end{tikzpicture}
\caption{The measured compression gap, reproduced on this machine. GPT-2 needs 5.23x
more tokens for identical text. At 0.779 chars/token it is emitting more than one token
per Urdu character, consistent with the two-bytes-per-character analysis in
Section~\ref{sec:whyurdu}.}
\label{fig:tokbar}
\end{figure}

\subsubsection{What the number means in money and context}

The claim is not an abstraction. Three concrete consequences follow. All three are
\textbf{arithmetic on the measured ratio}, not separate measurements, and are labelled
as such:

\begin{itemize}
\item \textbf{Training cost.} The v2 run processed 2.5B tokens for about \$13. The same
      text through GPT-2's tokenizer would be roughly 13B tokens, so roughly \$68 of
      GPU time for the same corpus, before any consideration of quality.
\item \textbf{Context window.} 1024 tokens holds about 4,176 Urdu characters with this
      tokenizer, and about 798 with GPT-2's. The same model architecture would see about
      a fifth as much Urdu at once.
\item \textbf{Learnability.} A model whose tokens are mostly single bytes must spend its
      early capacity learning to reassemble characters from bytes, before it can learn
      any Urdu. An Urdu-native tokenizer hands it meaningful units on the first step.
\end{itemize}

\begin{honest}{The README's stated command does not reproduce the README's number}
The README documents the measurement as
\code{python -m eval.tokenizer\_compare --gpt2 <gpt2 tokenizer.json>}, with no
\code{--max-chars}. But \code{--max-chars} defaults to \textbf{500,000}, and running the
command exactly as written yields a 504,251-character sample, 125,634 of our tokens
against GPT-2's 645,497, and a gain of \textbf{5.14x}, not 5.23x. Reproducing the
README's exact figures requires \code{--max-chars 800000}, which is the flag that
produces the cited 806,495-character sample.

Both numbers are real and neither is inflated (5.14x and 5.23x are the same finding).
The defect is purely documentary: the command printed next to the table is not the
command that produced the table. One flag added to the README closes it. Worth fixing
precisely because the number is genuine and reproduces perfectly once the flag is there;
it would be a shame to have a reviewer type the documented command, get a different
number, and start doubting the rest.
\end{honest}

\begin{honest}{The comparison samples from val, then train, which is fine but undocumented}
\code{sample\_text} iterates \code{for split in ["val", "train"]}. Since \file{data/val}
alone holds 8.7MB of text, the 800k-character sample never reaches \file{data/train} in
practice, so the sample is entirely held-out data. That is the right thing and it is what
"held-out Urdu sample" in the README means. But it is a consequence of the val directory
happening to be large enough, not an enforced property: with a smaller val split the
sample would silently start including training text, and the tokenizer would be
measured partly on data it was trained on. A single assertion, or dropping \code{"train"}
from that list, would make the guarantee real rather than incidental.
\end{honest}

\subsubsection{The byte-level fallback path}

If \code{--gpt2} is not supplied, the script compares against a synthetic worst case:
one token per UTF-8 byte. The docstring is careful to call this "the worst case GPT-2
approaches on unseen scripts" rather than GPT-2 itself. The measurement above shows how
close that approximation is: GPT-2's real 0.779 chars/token against the byte baseline's
roughly 0.5 for two-byte characters. GPT-2 does slightly better than pure bytes, because
a few merges do fire, but it is in the same regime. The fallback is an honest stand-in,
labelled as one.

\newpage
% =====================================================================
\section{Part V: The model}
\label{sec:model}

\file{model/transformer.py} is 186 lines and contains the entire network. Nothing is
imported from the \code{transformers} library. It is a modern decoder-only stack: not a
GPT-2 reimplementation, but the Llama-family design (RMSNorm, RoPE, SwiGLU, no biases,
tied embeddings).

\subsection{\code{config.py}: one dataclass, four presets}
\label{sec:presets}

\code{ModelConfig} is a \code{@dataclass} holding every architectural knob.
\code{\_\_post\_init\_\_} does two jobs: it expands the \code{n\_kv\_head == 0} sentinel
to mean "no grouped-query attention", and it asserts the two divisibility constraints the
rest of the code depends on (\code{n\_embd} divisible by \code{n\_head},
\code{n\_head} divisible by \code{n\_kv\_head}). Failing loudly at construction beats
failing with a shape mismatch fifty lines deep in attention.

\begin{keyidea}{Verified: every preset's parameter count matches the README}
The four presets were instantiated at the real 32k vocabulary and counted during the
writing of this document. All four match the README's table exactly.
\end{keyidea}

\begin{table}[H]
\centering
\small
\begin{tabular}{lrrrrrrr}
\toprule
\textbf{preset} & \textbf{layers} & \textbf{d\_model} & \textbf{heads} & \textbf{ctx}
& \textbf{embedding} & \textbf{blocks} & \textbf{total} \\
\midrule
tiny   &  6 &  256 &  4 &  512 &  8.19M &   4.82M &  \textbf{13.01M} \\
small  &  6 &  512 &  8 & 1024 & 16.38M &  19.27M &  \textbf{35.66M} \\
base   & 14 &  768 & 12 & 1024 & 24.58M &  99.11M & \textbf{123.69M} \\
medium & 24 & 1024 & 16 & 1024 & 32.77M & 303.61M & \textbf{336.38M} \\
\bottomrule
\end{tabular}
\caption{The four presets, measured on this machine, not copied from the README. The
embedding column is \code{vocab\_size} $\times$ \code{n\_embd}; because embeddings are
tied it is counted once, not twice.}
\label{tab:presets}
\end{table}

\begin{figure}[H]
\centering
\begin{tikzpicture}
\begin{axis}[
  width=13cm, height=5.4cm,
  xbar stacked, y=8mm, bar width=6mm,
  xmin=0, xmax=360,
  xlabel={parameters (millions)},
  symbolic y coords={medium,base,small,tiny},
  ytick=data, tick label style={font=\small}, label style={font=\small},
  legend style={font=\footnotesize, at={(0.985,0.95)}, anchor=north east, draw=linegrey,
    fill=white, fill opacity=0.92, text opacity=1},
  xmajorgrids, grid style={linegrey}, enlarge y limits=0.28,
]
\addplot[fill=warnred!45, draw=warnred] coordinates {(8.19,tiny) (16.38,small) (24.58,base) (32.77,medium)};
\addplot[fill=accent!55, draw=inkblue] coordinates {(4.82,tiny) (19.27,small) (99.11,base) (303.61,medium)};
\legend{embedding table, transformer blocks}
\end{axis}
\end{tikzpicture}
\caption{Where the parameters live. For \code{tiny}, the vocabulary table is 63\% of the
whole model; by \code{medium} it is under 10\%. This is the point the README makes
honestly rather than hiding.}
\label{fig:paramsplit}
\end{figure}

\begin{keyidea}{The "10M-parameter model" is a fiction, and the project says so}
The \code{tiny} preset's comment in \file{configs/tiny.yaml} still says "\textasciitilde{}10M params",
but a 32,000-token vocabulary at \code{d\_model}=256 costs 8.19M parameters before a
single transformer block exists. The blocks add only 4.82M. So a genuinely tiny
transformer still lands at 13.0M, and the initial parameter-count test failed at 15.2M
while the author was tuning it.

The choice made was to tune the block dimensions to a defensible 13M and \emph{document}
that 8.2M of it is vocabulary, rather than shrink the vocabulary to hit a round headline
number. That is the right call: shrinking the vocab to reach "10M" would have degraded
the very tokenizer the project exists to demonstrate, in exchange for a prettier number.
This is a good thing to be able to say out loud in an interview.
\end{keyidea}

\subsection{The block, drawn}

\begin{figure}[H]
\centering
\begin{tikzpicture}[node distance=5mm]
  % ==== main vertical spine ====
  \node[store,text width=3.0cm] (in) {token ids\\ \scriptsize (B, T) int64};
  \node[box,below=of in,text width=3.0cm] (emb) {\code{tok\_emb}\\ \scriptsize Embedding, 32000 x d\_model};
  \node[box,below=of emb,text width=3.0cm] (drop) {\code{drop}\\ \scriptsize Dropout};

  \node[box,below=8mm of drop,text width=3.0cm,fill=softgrey] (n1) {\code{attn\_norm}\\ \scriptsize RMSNorm};
  \node[box,below=of n1,text width=3.0cm] (att) {\code{attn}\\ \scriptsize RoPE causal self-attention};
  \node[dec,below=of att] (add1) {$+$};
  \node[box,below=of add1,text width=3.0cm,fill=softgrey] (n2) {\code{mlp\_norm}\\ \scriptsize RMSNorm};
  \node[box,below=of n2,text width=3.0cm] (mlp) {\code{mlp}\\ \scriptsize SwiGLU};
  \node[dec,below=of mlp] (add2) {$+$};

  \node[box,below=8mm of add2,text width=3.0cm,fill=softgrey] (fn) {\code{final\_norm}\\ \scriptsize RMSNorm};
  \node[box,below=of fn,text width=3.0cm] (head) {\code{lm\_head}\\ \scriptsize Linear, no bias};
  \node[store,below=of head,text width=3.0cm] (out) {logits\\ \scriptsize (B, T, 32000)};

  \draw[flow] (in) -- (emb);
  \draw[flow] (emb) -- (drop);
  \draw[flow] (drop) -- (n1);
  \draw[flow] (n1) -- (att);
  \draw[flow] (att) -- (add1);
  \draw[flow] (add1) -- (n2);
  \draw[flow] (n2) -- (mlp);
  \draw[flow] (mlp) -- (add2);
  \draw[flow] (add2) -- (fn);
  \draw[flow] (fn) -- (head);
  \draw[flow] (head) -- (out);

  % ==== residual skips, down the left ====
  \draw[dflow] (drop.west) to[out=180,in=180,looseness=1.5] (add1.west);
  \draw[dflow] (add1.west) to[out=180,in=180,looseness=1.5] (add2.west);
  \node[lbl,left=1mm of $(drop.west)!0.5!(add1.west)$,text width=2.1cm]
    {residual:\\ attention writes\\ a correction};
  \node[lbl,left=1mm of $(add1.west)!0.5!(add2.west)$,text width=2.1cm]
    {residual:\\ the MLP writes\\ a correction};

  % ==== RoPE cache feeds attention from the right ====
  \node[extbox,right=9mm of att,text width=2.7cm] (rope) {\code{\_rope(...)}\\ \scriptsize cos/sin cache,\\ applied to q and k only};
  \draw[dflow] (rope) -- (att);

  % ==== tied weights: routed well clear on the far right ====
  \draw[dflow] (emb.east) -- ++(6.4,0) coordinate (tie) |- (head.east);
  \node[lbl,anchor=west,text width=2.7cm] at ($(tie)+(0.15,-3.1)$)
    {tied: one weight matrix,\\ used twice\\ \scriptsize (saves 24.6M params\\ in \code{base})};

  \begin{scope}[on background layer]
    \node[draw=accent,dashed,rounded corners=3pt,inner sep=6pt,
          fit=(n1)(att)(add1)(n2)(mlp)(add2)] (blk) {};
  \end{scope}
  \node[lbl,anchor=west,text width=2.5cm] at ($(blk.south east)+(0.3,0.75)$)
    {\code{Block},\\ repeated \code{n\_layer}\\ times (14 for \code{base})};
\end{tikzpicture}
\caption{The full forward pass. Dashed arrows are residual connections and the tied
weight matrix. Note that both norms sit \emph{before} their sublayer, not after.}
\label{fig:block}
\end{figure}

\begin{jargon}{Residual connection}
The \code{x = x + sublayer(norm(x))} pattern. The sublayer computes a \emph{correction}
that is added to its input, rather than replacing it. This gives gradients a clean path
straight from the loss back to the earliest layers, which is what makes deep stacks
trainable at all. Both \code{+} nodes in Figure~\ref{fig:block} are these.
\end{jargon}

\begin{jargon}{Pre-norm vs post-norm}
Where the normalization sits. The original transformer normalized \emph{after} the
sublayer (post-norm), which needs careful learning-rate warmup to avoid diverging.
Pre-norm, used here, normalizes the sublayer's input and leaves the residual path
untouched. It is markedly more stable, and it is why the residual stream in
Figure~\ref{fig:block} runs from embedding to \code{final\_norm} without ever passing
through a normalization.
\end{jargon}

\subsection{\code{RMSNorm}}

\begin{lstlisting}[language=Python,caption={RMSNorm in full}]
class RMSNorm(nn.Module):
    def __init__(self, dim, eps=1e-5):
        super().__init__()
        self.eps = eps
        self.weight = nn.Parameter(torch.ones(dim))

    def forward(self, x):
        dtype = x.dtype
        x = x.float()                                    # fp32 for stability
        x = x * torch.rsqrt(x.pow(2).mean(-1, keepdim=True) + self.eps)
        return (x * self.weight.float()).to(dtype)
\end{lstlisting}

\begin{jargon}{Normalization, and RMSNorm specifically}
As numbers flow through layers their scale drifts, and drift destabilizes training.
Normalization rescales each vector to a standard size. The older LayerNorm subtracts the
mean and divides by the standard deviation, then applies a learned scale and shift.
RMSNorm drops the mean-subtraction and the shift entirely: it divides by the root mean
square and applies one learned scale. It works about as well, with fewer operations and
fewer parameters, which is why every recent model uses it.
\end{jargon}

Two details in those five lines are load-bearing:

\begin{itemize}
\item \code{eps} inside the square root prevents division by zero for an all-zero vector.
\item The \code{x.float()} / \code{.to(dtype)} sandwich is the important one. Training
      runs in bfloat16 (Section~\ref{sec:amp}), which has very few mantissa bits.
      Squaring and averaging in bf16 loses precision exactly where it matters, so the
      statistic is computed in fp32 and only the result is cast back. This is a standard,
      deliberate mixed-precision idiom, not an accident.
\end{itemize}

\subsection{Rotary position embeddings (RoPE)}

\begin{jargon}{Positional encoding, and why it is needed}
Attention is a weighted average, and averages have no order: without help, "dog bites
man" and "man bites dog" look identical to it. Something must inject position. GPT-2
added a learned position vector to each token's embedding. RoPE takes a different route.
\end{jargon}

\begin{jargon}{RoPE}
Instead of \emph{adding} position information, RoPE \emph{rotates} each query and key
vector by an angle proportional to its position. Pairs of dimensions are treated as
points on a plane and spun. The elegance is what happens next: when attention takes the
dot product of a query at position $m$ with a key at position $n$, the rotations combine
so the result depends only on the \emph{difference} $m-n$. Relative position falls out of
the arithmetic for free, with zero parameters.
\end{jargon}

\begin{lstlisting}[language=Python,caption={Building and applying the rotation}]
def build_rope_cache(seq_len, head_dim, theta, device, dtype):
    inv_freq = 1.0 / (theta ** (torch.arange(0, head_dim, 2).float() / head_dim))
    t = torch.arange(seq_len).float()
    freqs = torch.outer(t, inv_freq)      # (seq_len, head_dim/2)
    return freqs.cos().to(dtype), freqs.sin().to(dtype)

def apply_rope(x, cos, sin):              # x: (B, n_head, T, head_dim)
    T = x.shape[-2]
    cos = cos[:T].unsqueeze(0).unsqueeze(0)
    sin = sin[:T].unsqueeze(0).unsqueeze(0)
    x1, x2 = x.chunk(2, dim=-1)
    rot = torch.cat((-x2, x1), dim=-1)    # the 90-degree rotation
    cos_full = torch.cat((cos, cos), dim=-1)
    sin_full = torch.cat((sin, sin), dim=-1)
    return x * cos_full + rot * sin_full
\end{lstlisting}

\code{inv\_freq} builds a spectrum of rotation speeds from \code{theta} (10,000 here):
the first dimension pair rotates fast and distinguishes neighbours, the last rotates
extremely slowly and distinguishes distant positions. Together they encode position at
every scale. \code{rot = cat((-x2, x1))} is the standard trick for rotating every pair at
once: it is what $(a,b) \rightarrow (-b,a)$ does, vectorized across the whole tensor.

\begin{keyidea}{RoPE is applied to queries and keys only, never to values}
\code{q = apply\_rope(q, cos, sin)} and \code{k = apply\_rope(k, cos, sin)}, but
\code{v} is untouched. This is correct and it is the whole design: position must
influence \emph{who attends to whom} (which is decided by the query-key dot product),
but the content that gets passed along once a position is selected should not be
scrambled by where it happened to sit. Getting this wrong is a classic RoPE bug; this
code gets it right.
\end{keyidea}

\begin{honest}{The RoPE cache key is subtly wasteful}
\code{\_rope(self, T, device, dtype)} keys its cache on \code{(T, device, dtype)}, but
the value it builds ignores \code{T} entirely and always constructs the cache at
\code{self.cfg.max\_seq\_len}. So a forward pass at \code{T=64} and one at \code{T=128}
create two cache entries holding \textbf{identical} tensors. Nothing is incorrect
(\code{apply\_rope} slices \code{cos[:T]}, so any entry long enough works), and training
always uses one fixed \code{T} so exactly one entry is ever built. But \code{generate()}
calls the model with a growing \code{T} on every single step, so sampling 60 tokens
builds up to 60 identical full-size caches in a plain dict that is never evicted. For the
\code{base} model that is 60 copies of a 1024x32 tensor: small in absolute terms, but
pointless. Keying on \code{(device, dtype)} alone would be strictly better and is a
one-line change.

A related note: \code{\_rope\_cache} is a plain Python dict, not a registered buffer, so
it is invisible to \code{.to()}, \code{state\_dict()} and \code{.cuda()}. That is
actually \emph{why} the device is in the key, and it means checkpoints do not carry
redundant cache tensors. Defensible, but it is the kind of thing that looks like an
oversight until you work out why it is there.
\end{honest}

\subsection{\code{Attention}}

\begin{lstlisting}[language=Python,caption={The attention forward pass}]
q = self.q_proj(x).view(B, T, self.n_head, self.head_dim).transpose(1, 2)
k = self.k_proj(x).view(B, T, self.n_kv_head, self.head_dim).transpose(1, 2)
v = self.v_proj(x).view(B, T, self.n_kv_head, self.head_dim).transpose(1, 2)

q = apply_rope(q, cos, sin)
k = apply_rope(k, cos, sin)

if self.n_kv_head != self.n_head:              # grouped-query attention
    rep = self.n_head // self.n_kv_head
    k = k.repeat_interleave(rep, dim=1)
    v = v.repeat_interleave(rep, dim=1)

y = F.scaled_dot_product_attention(q, k, v, is_causal=True,
                                   dropout_p=self.dropout if self.training else 0.0)
y = y.transpose(1, 2).contiguous().view(B, T, self.n_head * self.head_dim)
return self.o_proj(y)
\end{lstlisting}

\begin{jargon}{Multi-head attention}
Rather than one attention computation over 768 dimensions, run 12 independent ones over
64 dimensions each and concatenate. Each "head" can specialize: one might track
subject-verb agreement, another long-range topic. The \code{view(...).transpose(1,2)}
pair is what reshapes \code{(B, T, 768)} into \code{(B, 12, T, 64)} so the heads compute
in parallel as a batch dimension.
\end{jargon}

\begin{jargon}{Grouped-query attention (GQA)}
An inference optimization. Generation caches every past key and value, and that cache
dominates memory. GQA gives each head its own query but shares keys and values across
groups of heads, shrinking the cache proportionally. Here \code{n\_kv\_head} controls it,
and \code{repeat\_interleave} expands the shared keys back out to per-head shape.
\end{jargon}

\begin{honest}{GQA is implemented but not used by any preset, and its implementation is the naive one}
Every one of the four presets leaves \code{n\_kv\_head} at its default 0, which
\code{\_\_post\_init\_\_} expands to \code{n\_head}, which makes
\code{n\_kv\_head != n\_head} false and skips the branch entirely. So no shipped model
uses GQA; the README's "grouped-query-capable" is exactly the right word, and the code is
honest about it.

Two things are worth knowing anyway. First, the \code{repeat\_interleave} implementation
materializes the expanded keys and values in full, which means it saves \emph{no memory
at all} during the forward pass. It is a correctness shim, not an optimization. Real GQA
savings come from the KV cache during generation, and this model has no KV cache (see
below), so there is currently nothing for GQA to save. Second, that is fine: the branch
is there so the config surface is honest and a future run can turn it on. It just should
not be described as a performance feature, because as implemented it is not one.
\end{honest}

\begin{keyidea}{One call does the mask, the scaling and the Flash kernel}
\code{F.scaled\_dot\_product\_attention(q, k, v, is\_causal=True)} is doing four things
that older implementations spell out by hand: it scales by $1/\sqrt{d_{head}}$, builds
and applies the causal mask, softmaxes, and dispatches to a fused FlashAttention kernel
on CUDA when the shapes allow. \code{is\_causal=True} means the triangular mask in
Figure~\ref{fig:causal} is never materialized as a tensor at all, which is where a large
part of the memory saving comes from. The README's claim that "RoPE plus
\code{scaled\_dot\_product\_attention(is\_causal=True)} is enough to get a correct,
Flash-ready causal transformer with no hand-written mask" is accurate.
\end{keyidea}

\begin{jargon}{FlashAttention}
An attention implementation that never writes the full (T x T) attention matrix to GPU
memory, computing it in tiles inside fast on-chip memory instead. It is exact, not an
approximation, and it is the reason a 1024-token context is cheap. It is what delivered
the measured $\sim$105,000 tokens/second in Section~\ref{sec:runs}.
\end{jargon}

\code{bias=False} on all four projections is another Llama-family choice: biases
contribute little in transformers and cost parameters and a kernel launch.
\code{contiguous()} before the final \code{view} is required because \code{transpose}
returns a non-contiguous tensor that \code{view} cannot reshape.

\subsection{\code{SwiGLU}}

\begin{lstlisting}[language=Python,caption={The MLP}]
hidden = int(cfg.mlp_ratio * cfg.n_embd)     # mlp_ratio = 8/3
hidden = 64 * ((hidden + 63) // 64)          # round up to a multiple of 64
self.gate_proj = nn.Linear(cfg.n_embd, hidden, bias=False)
self.up_proj   = nn.Linear(cfg.n_embd, hidden, bias=False)
self.down_proj = nn.Linear(hidden, cfg.n_embd, bias=False)

def forward(self, x):
    return self.down_proj(F.silu(self.gate_proj(x)) * self.up_proj(x))
\end{lstlisting}

\begin{jargon}{MLP / feed-forward block}
Attention moves information between positions; the MLP does the thinking at each
position independently. Classically it projects up to 4x the width, applies a non-linear
function, and projects back down.
\end{jargon}

\begin{jargon}{SwiGLU, and the gating idea}
A gated variant. The input is projected \emph{twice}, into a "gate" and an "up" path. The
gate goes through SiLU ($x \cdot \sigma(x)$, a smooth version of ReLU) and then
\textbf{multiplies} the up path elementwise. So the gate learns, per dimension, how much
of the up path to let through. That multiplicative interaction is more expressive than a
single projection, and SwiGLU consistently outperforms plain GELU at equal parameter
count.
\end{jargon}

\begin{keyidea}{Why the ratio is 8/3, and why base needs 14 layers rather than 12}
SwiGLU needs \textbf{three} weight matrices where a classic MLP needs two. At the
textbook 4x width that would be 50\% more parameters than a GELU block. The 8/3 ratio
exists to cancel exactly that: $3 \times \frac{8}{3} = 8 = 2 \times 4$, so a SwiGLU block
at 8/3 has the same parameter count as a GELU block at 4x.

The consequence is the one the README calls out. Each \code{base} layer is
parameter-cheaper than a GPT-2 layer of the same width, so the textbook 12 layers do not
reach 124M. The preset uses \textbf{14}. Depth, not the memorized layer count, is what
hits a target size. This is a genuinely non-obvious thing to have worked out and it is a
strong interview answer.
\end{keyidea}

The rounding to a multiple of 64 is a GPU alignment concern: tensor-core matrix
multiplies are fastest on dimensions that are multiples of 64, so the ratio is treated as
a target rather than an exact requirement.

\subsection{\code{UrduLM}: assembly, tying, and initialization}

\subsubsection{Weight tying}

\begin{lstlisting}[language=Python,caption={Tying, in two lines}]
self.lm_head = nn.Linear(cfg.n_embd, cfg.vocab_size, bias=False)
if cfg.tie_embeddings:
    self.lm_head.weight = self.tok_emb.weight
\end{lstlisting}

\begin{jargon}{Weight tying}
Using one matrix for both the input embedding (id to vector) and the output projection
(vector to scores over ids). They are natural transposes of each other: both encode "what
does token $i$ look like in vector space". Sharing them saves \code{vocab\_size} $\times$
\code{n\_embd} parameters (24.6M for \code{base}) and usually improves quality, because
one matrix gets gradient signal from both directions.
\end{jargon}

This is a genuine assignment of the same Python object, not a copy, so the two stay
identical for the whole of training. \code{tests/test\_model.py} asserts it at the level
that actually proves it: \code{model.lm\_head.weight.data\_ptr() == model.tok\_emb.weight.data\_ptr()},
comparing the memory addresses rather than the values. Comparing values would pass even
if they were two matrices that merely started equal and then diverged. That is a
well-designed test.

\subsubsection{Initialization}

\begin{lstlisting}[language=Python,caption={The two-stage init}]
self.apply(self._init_weights)                    # everything: normal(0, 0.02)
for name, p in self.named_parameters():           # then override the residual outputs
    if name.endswith("o_proj.weight") or name.endswith("down_proj.weight"):
        nn.init.normal_(p, mean=0.0, std=0.02 / math.sqrt(2 * cfg.n_layer))
\end{lstlisting}

\begin{keyidea}{Why the residual projections get a smaller initial scale}
Every block \emph{adds} to the residual stream. With $L$ blocks each contributing
independent noise, the stream's variance grows roughly linearly with depth, so a deep
stack starts with a wildly-scaled output before it has learned anything. Scaling the two
projections that write into the residual stream (\code{o\_proj} from attention,
\code{down\_proj} from the MLP) by $1/\sqrt{2L}$ keeps the accumulated variance roughly
constant regardless of depth. The factor 2 is because each layer writes twice, once from
attention and once from the MLP. This is the GPT-2 initialization scheme, and it is
correctly reproduced here.
\end{keyidea}

\subsubsection{\code{forward}}

The forward pass asserts \code{T <= cfg.max\_seq\_len} (a real guard: RoPE's cache is
built to exactly that length), embeds, runs the blocks, normalizes, projects to logits,
and optionally computes the loss with \code{ignore\_index=-100}.

\begin{jargon}{\code{ignore\_index=-100}}
PyTorch's convention for "do not score this position". Any target set to $-100$
contributes nothing to the loss. Pretraining never uses it (every position is a real
target), but \code{eval/cloze.py} depends on it entirely: it masks out the context so
only the candidate answer's tokens are scored. See Section~\ref{sec:cloze}.
\end{jargon}

\subsubsection{\code{generate}}

\begin{lstlisting}[language=Python,caption={Autoregressive sampling}]
for _ in range(max_new_tokens):
    idx_cond = idx[:, -self.cfg.max_seq_len:]        # crop to the context window
    logits, _ = self(idx_cond)
    logits = logits[:, -1, :] / max(temperature, 1e-6)
    if top_k is not None:
        v, _ = torch.topk(logits, min(top_k, logits.size(-1)))
        logits[logits < v[:, [-1]]] = -float("inf")
    probs = F.softmax(logits, dim=-1)
    next_id = torch.multinomial(probs, num_samples=1)
    idx = torch.cat((idx, next_id), dim=1)
\end{lstlisting}

\begin{jargon}{Temperature}
A divisor on the logits before softmax. Below 1 it sharpens the distribution (more
predictable, more repetitive); above 1 it flattens it (more varied, more likely to be
nonsense). \code{max(temperature, 1e-6)} guards against a division by zero when someone
asks for greedy decoding by passing 0.
\end{jargon}

\begin{jargon}{Top-k sampling}
Keep only the $k$ highest-scoring tokens and set the rest to negative infinity, which
softmax turns into exactly zero probability. This prevents the long tail of thousands of
barely-plausible tokens from collectively stealing a meaningful share of the draws. The
project samples at \code{top\_k=40}.
\end{jargon}

\begin{honest}{\code{generate} has no KV cache, so sampling is quadratic}
Every iteration re-runs the \emph{entire} forward pass over the whole prompt so far, and
then throws away all of it except the last position's logits. Generating $n$ tokens
therefore costs $O(n^2)$ work instead of the $O(n)$ that caching past keys and values
would give. Standard implementations keep a KV cache exactly for this.

For this project's actual usage (a handful of 60 to 80 token samples for the script-mix
diagnostic) the waste is invisible and the simplicity is worth more. But it means this
model is not deployable for interactive serving as written, and it is the reason the GQA
branch above has nothing to optimize: GQA shrinks a KV cache, and there is no KV cache.
Naming this before an interviewer does is much better than the reverse.
\end{honest}

\begin{honest}{\code{num\_params(non\_embedding=True)} has two branches that do the same thing}
\begin{lstlisting}[language=Python]
n = sum(p.numel() for p in self.parameters())
if non_embedding and not self.cfg.tie_embeddings:
    n -= self.tok_emb.weight.numel()
if non_embedding and self.cfg.tie_embeddings:
    n -= self.tok_emb.weight.numel()
\end{lstlisting}
The two conditions are exhaustive and mutually exclusive, and their bodies are
identical, so the whole thing is equivalent to a single
\code{if non\_embedding:}. It reads as though it were meant to handle the two cases
differently (in the untied case \code{sum(parameters())} counts both the embedding and
\code{lm\_head}, so subtracting only \code{tok\_emb} still leaves the head's copy in the
total) and then collapsed. The result is not wrong for any shipped preset, since all four
tie embeddings and nothing calls the flag anyway. It is dead branching that should be one
line, and if the untied case is ever used the "non-embedding" count will silently include
the output head.
\end{honest}

\subsection{The tests}
\label{sec:tests}

\begin{keyidea}{Verified: all 13 tests pass on this machine}
\code{python -m pytest tests/ -q} was run during the writing of this document:
\textbf{13 passed in 55.42s}, matching the README's claim of 13 tests (5 model, 8
pipeline).
\end{keyidea}

The best of them is \code{test\_causal\_masking}, because it tests the \emph{property}
rather than the implementation:

\begin{lstlisting}[language=Python,caption={Testing causality by experiment}]
base, _ = model(x)
x2 = x.clone()
x2[0, -1] = (x2[0, -1] + 1) % cfg.vocab_size     # perturb the LAST token
pert, _ = model(x2)
assert torch.allclose(base[:, :-1], pert[:, :-1], atol=1e-5)      # earlier: unchanged
assert not torch.allclose(base[:, -1], pert[:, -1], atol=1e-5)    # last: changed
\end{lstlisting}

Change a future token; if any earlier position's output moves, information flowed
backwards and the mask is broken. The second assertion is what makes it a real test
rather than a tautology: without it, a model that ignored its input entirely would pass.
This would catch a swapped mask, a missing \code{is\_causal}, or an off-by-one in RoPE,
none of which a shape test would notice.

\code{test\_preset\_param\_counts} asserts each preset lands inside a band (for instance
\code{base} in 115M to 135M). Bands rather than exact numbers is the right choice: it
pins the intent without making every dimension tweak a test failure.

\newpage
% =====================================================================
\section{Part VI: Training}
\label{sec:training}

\file{train.py} is 268 lines and is the piece that had to work unattended on a metered
GPU. Its design priorities show that: budget in tokens rather than steps, resume
exactly, log everything, and never surprise the operator.

\subsection{\code{TokenLoader}: batches from a memory-mapped stream}

\begin{lstlisting}[language=Python,caption={Setting up the sampler}]
self.data = np.memmap(bin_path, dtype=np.uint16, mode="r")
n_blocks = (len(self.data) - 1) // block_size
self.starts = np.arange(n_blocks) * block_size
rng = np.random.default_rng(seed)
rng.shuffle(self.starts)          # fixed shuffle, seeded
self.cursor = 0
\end{lstlisting}

The whole corpus is one long array of token ids. The loader precomputes every valid block
start, shuffles that list once with a seeded generator, and walks it with a cursor.

\begin{keyidea}{Why the cursor is the entire resume mechanism}
The shuffle is deterministic given the seed, so \code{self.starts} is the \emph{same
list in the same order} on every run of every process. That means the loader's complete
state is one integer: how far down the list it has walked. \code{state\_dict()} returns
\code{\{"cursor": self.cursor\}} and nothing else. Restore that integer and the exact
same batch sequence resumes. This is a genuinely elegant reduction, and it is why resume
here is trustworthy rather than approximate.
\end{keyidea}

\begin{lstlisting}[language=Python,caption={Drawing one batch}]
idxs = []
need = self.batch_size * self.world_size
while len(idxs) < need:
    if self.cursor >= len(self.starts):
        self.cursor = 0                      # epoch wrap
    idxs.append(self.starts[self.cursor])
    self.cursor += 1
idxs = idxs[self.rank::self.world_size][: self.batch_size]
x = np.stack([self.data[i : i + self.block_size].astype(np.int64) for i in idxs])
y = np.stack([self.data[i + 1 : i + 1 + self.block_size].astype(np.int64) for i in idxs])
\end{lstlisting}

\begin{keyidea}{\code{x} and \code{y} are the same slice, offset by one}
This is the whole of next-token prediction, in two lines. \code{y} is \code{x} shifted
left by one position, so the target at every position is simply the next token. Combined
with causal masking, one forward pass over a 1024-token block trains 1024 separate
predictions at once. That density is why pretraining works at all.
\end{keyidea}

Every rank draws the \emph{same} \code{need} indices and then takes its own stride
(\code{idxs[rank::world\_size]}). This is deliberate: since all ranks advance the cursor
identically, the cursor stays synchronized across processes without any communication,
and rank 0's saved cursor is valid for the whole job.

\begin{honest}{Three real limitations of the loader}
\begin{itemize}
\item \textbf{The epoch wrap does not reshuffle.} On \code{cursor = 0} the same order
      repeats. The v2 run made roughly 3 passes over its data, so it saw the identical
      block sequence three times. Reshuffling with a per-epoch seed is standard and is
      strictly better; the cost of not doing it is small but real.
\item \textbf{Blocks are cut at fixed stride offsets.} Starts are always multiples of
      \code{block\_size}, so the corpus is carved into a fixed tiling and no block ever
      straddles a boundary differently between epochs. A random-offset scheme would give
      more distinct contexts from the same data. Again small, again free.
\item \textbf{No \code{DataLoader}, no prefetch thread.} Batches are built synchronously
      in the training loop, so the GPU idles during the numpy slicing and the host-to-
      device copy. At the measured $\sim$105,000 tokens/second this evidently was not the
      bottleneck, so the simplicity was the right trade. It is worth knowing it is a
      trade.
\end{itemize}
\end{honest}

\begin{honest}{Documents are not respected: blocks straddle \code{<|endoftext|>} freely}
Nothing aligns block boundaries to document boundaries. A 1024-token block routinely
contains the tail of one document, an \code{<|endoftext|>}, and the head of the next, and
the model is asked to predict across that seam using unrelated context. This is exactly
what GPT-2 and nanoGPT do and is entirely conventional, so it is not a defect. It is
worth being able to say why it is acceptable: the \code{<|endoftext|>} token is the
signal that teaches the model a document ended, and attention learns to discount context
before it.
\end{honest}

\subsection{The optimizer and the schedule}

\begin{jargon}{Gradient descent, and AdamW}
Training computes, for every parameter, the direction that would reduce the loss (the
\emph{gradient}), and steps a little way along it. Plain descent uses the raw gradient.
Adam keeps running averages of each parameter's recent gradients and their squared
magnitudes, and uses them to give every parameter its own effective step size, which is
far more robust. AdamW adds \emph{decoupled weight decay}: a gentle pull of every weight
toward zero, applied separately from the gradient step, which regularizes without
corrupting Adam's statistics.
\end{jargon}

The configuration is \code{betas=(0.9, 0.95)} and \code{weight\_decay=0.1}. Both are the
GPT-3 paper's values and are conventional for this model class: \code{beta2=0.95} rather
than the PyTorch default 0.999 makes the second-moment estimate adapt faster, which is
the standard choice for language models.

\begin{lstlisting}[language=Python,caption={The learning-rate schedule, in seven lines}]
def cosine_lr(step, warmup, total, base_lr, min_lr):
    if step < warmup:
        return base_lr * (step + 1) / warmup      # linear warmup
    if step >= total:
        return min_lr
    ratio = (step - warmup) / max(1, total - warmup)
    coeff = 0.5 * (1 + math.cos(math.pi * ratio))  # 1 -> 0 over the run
    return min_lr + coeff * (base_lr - min_lr)
\end{lstlisting}

\begin{jargon}{Learning rate, warmup, and cosine decay}
The learning rate is how far each step moves. It is not constant. \emph{Warmup} ramps it
up linearly from near zero over the first steps, because a randomly initialized model
produces enormous unreliable gradients and a full-size step on them can destroy the model
in the first minute. \emph{Cosine decay} then eases it smoothly down to a small floor, so
early training explores broadly and late training settles precisely.
\end{jargon}

\begin{figure}[H]
\centering
\begin{tikzpicture}
\begin{axis}[
  width=13.2cm, height=5.2cm,
  xlabel={step}, ylabel={learning rate},
  xmin=0, xmax=4768, ymin=0, ymax=0.00033,
  tick label style={font=\small}, label style={font=\small},
  xmajorgrids, ymajorgrids, grid style={linegrey},
  scaled y ticks=false, yticklabel style={/pgf/number format/fixed,
     /pgf/number format/precision=5},
  legend style={font=\footnotesize, at={(0.98,0.95)}, anchor=north east, draw=linegrey},
]
\addplot[draw=accent, very thick, samples=300, domain=0:1000] {0.0003*(x+1)/1000};
\addplot[draw=inkblue, very thick, samples=300, domain=1000:4768]
  {0.00003 + 0.5*(1+cos(deg(pi*(x-1000)/3768)))*(0.0003-0.00003)};
\legend{linear warmup (1000 steps), cosine decay to min\_lr}
\end{axis}
\end{tikzpicture}
\caption{The schedule \code{cosine\_lr} produces for the \code{base} preset's committed
values (\code{lr} 3e-4, \code{min\_lr} 3e-5, \code{warmup\_steps} 1000), drawn over the
v2 run's actual 4,768 steps.}
\label{fig:lr}
\end{figure}

\begin{honest}{Warmup divides by \code{warmup}, not \code{warmup-1}, so the peak LR is never reached}
\code{base\_lr * (step + 1) / warmup} at \code{step = warmup - 1} gives exactly
\code{base\_lr}, and then the next step takes the cosine branch at \code{ratio = 0},
which also gives \code{base\_lr}. So the schedule is continuous and correct. This is worth
checking rather than assuming, because an off-by-one here is a classic silent bug that
either skips the peak or double-counts it. This code gets it right.
\end{honest}

\subsection{Token-count budgeting}

\begin{lstlisting}[language=Python]
tokens_per_step = batch_size * block_size * grad_accum * world_size
total_steps = max(1, max_tokens // tokens_per_step)
warmup_steps = cfg.get("warmup_steps", max(1, int(0.03 * total_steps)))
\end{lstlisting}

\begin{keyidea}{Budget in tokens, derive steps}
The config names \code{max\_tokens}, not \code{max\_steps}. Steps are derived. This is
the right primitive: "step" is meaningless without knowing the batch size, so a
step-budgeted config silently changes the amount of training whenever the batch changes,
while a token budget stays comparable across hardware, batch sizes and world sizes. It is
also what makes the Chinchilla reasoning in \file{docs/GPU\_RUN.md} directly expressible
in the config.

For the \code{base} preset: $32 \times 1024 \times 16 \times 1 = 524{,}288$ tokens per
step. This is verifiable against the committed log: the v2 run's records step in
increments of exactly 10,485,760 tokens per 20 steps, which is $20 \times 524{,}288$.
\end{keyidea}

\subsection{Gradient accumulation, and mixed precision}
\label{sec:amp}

\begin{jargon}{Gradient accumulation}
A large batch is statistically desirable but may not fit in VRAM. Accumulation runs
\code{grad\_accum} smaller "micro-batches", each doing a backward pass that \emph{adds}
into the same gradient buffers, and only then takes one optimizer step. The result is
mathematically the same as one large batch, at the cost of more time. Here 16
micro-batches of 32 sequences give an effective batch of 512 sequences, or 524,288
tokens.
\end{jargon}

\begin{lstlisting}[language=Python,caption={The accumulation loop}]
optimizer.zero_grad(set_to_none=True)
loss_accum = 0.0
for micro in range(grad_accum):
    x, y = train_loader.next_batch()
    if is_dist():
        model.require_backward_grad_sync = (micro == grad_accum - 1)
    with ctx:
        _, loss = model(x, y)
        loss = loss / grad_accum          # <- so the sum is a mean
    loss.backward()
    loss_accum += loss.item()
torch.nn.utils.clip_grad_norm_(raw_model.parameters(), grad_clip)
optimizer.step()
\end{lstlisting}

Three correct details:

\begin{itemize}
\item \code{loss / grad\_accum} is essential. Backward passes \emph{sum} into the
      gradient buffers, so without the division the effective gradient would be 16x too
      large and the learning rate 16x too high.
\item \code{require\_backward\_grad\_sync = (micro == grad\_accum - 1)} is the DDP
      optimization that matters most. By default DDP synchronizes gradients across GPUs
      on \emph{every} backward pass; here it is suppressed for the first 15
      micro-batches and allowed only on the last, cutting inter-GPU communication by
      16x. Getting this right is the difference between DDP helping and DDP being a
      bottleneck.
\item \code{set\_to\_none=True} frees the gradient tensors rather than filling them with
      zeros: slightly less memory, slightly faster.
\end{itemize}

\begin{jargon}{Gradient clipping}
If a batch produces a freak gradient, one step can wreck the model. \code{clip\_grad\_norm\_}
measures the total magnitude of all gradients and, if it exceeds \code{grad\_clip} (1.0
here), scales them all down proportionally. It preserves the direction and caps the size.
This is cheap insurance for a long unattended run.
\end{jargon}

\begin{jargon}{bfloat16 and AMP}
\emph{Automatic Mixed Precision} runs most operations in a 16-bit format for speed and
memory while keeping the master weights in 32-bit. \emph{bfloat16} keeps float32's full
exponent range and sacrifices mantissa bits, which means it cannot overflow where fp32
would not, and therefore needs no loss scaling at all (unlike fp16). It is the reason
\code{torch.cuda.amp.GradScaler} appears nowhere in this file.
\end{jargon}

\begin{lstlisting}[language=Python,caption={The precision decision}]
use_bf16 = device == "cuda" and torch.cuda.is_bf16_supported()
ctx = (torch.autocast(device_type="cuda", dtype=torch.bfloat16)
       if use_bf16 else torch.autocast(device_type="cpu", enabled=False))
\end{lstlisting}

The fallback is an \code{autocast} context with \code{enabled=False}, which is a no-op
context manager. That is a tidy way to keep one \code{with ctx:} in the loop instead of
branching, and CPU training runs in plain fp32.

\subsection{Checkpointing and resume}

\begin{figure}[H]
\centering
\resizebox{0.95\textwidth}{!}{
\begin{tikzpicture}[node distance=7mm]
  \node[box,text width=2.9cm] (cfg) {read YAML config\\ \scriptsize CLI flags override};
  \node[box,right=of cfg,text width=2.6cm] (ddp) {\code{setup\_ddp()}\\ \scriptsize rank, world\_size};
  \node[box,right=of ddp,text width=2.7cm] (build) {build model\\ \scriptsize + AdamW + loaders};
  \node[dec,right=of build,text width=1.5cm] (r) {\code{--resume}\\ and ckpt?};
  \node[box,below=9mm of r,text width=3.4cm] (load) {restore model, optimizer,\\ cursor, step, tokens\_seen};
  \node[box,left=10mm of load,text width=2.4cm] (loop) {\textbf{training loop}\\ \scriptsize while step < total};
  \node[box,below=9mm of loop,text width=3.0cm] (log) {every 20 steps:\\ \scriptsize append metrics.jsonl};
  \node[box,left=9mm of log,text width=3.0cm] (evl) {every 1000 steps:\\ \scriptsize val loss + ppl};
  \node[box,below=9mm of log,text width=3.0cm] (ck) {every 2000 steps:\\ \scriptsize write ckpt.pt};
  \node[store,below=9mm of ck,text width=3.4cm] (fin) {at the end: ckpt.pt (fat)\\ + model.pt (slim)};

  \draw[flow] (cfg) -- (ddp);
  \draw[flow] (ddp) -- (build);
  \draw[flow] (build) -- (r);
  \draw[flow] (r) -- node[lbl,right] {yes} (load);
  \draw[flow] (load) -- (loop);
  \draw[flow] (r.north) to[out=90,in=90,looseness=1.1] node[lbl,above] {no, start fresh} (loop.north);
  \draw[flow] (loop) -- (log);
  \draw[flow] (log) -- (evl);
  \draw[flow] (log) -- (ck);
  \draw[flow] (ck) -- (fin);
  \node[extbox,text width=3.6cm,anchor=west] at ($(fin.east)+(0.4,0)$)
    {\scriptsize master rank (rank 0) only\\ writes logs and checkpoints};
\end{tikzpicture}}
\caption{The control flow of \file{train.py}. Intervals shown are the \code{base} preset's
committed values.}
\label{fig:trainflow}
\end{figure}

Two checkpoints are written, and the distinction is deliberate:

\begin{table}[H]
\centering
\small
\begin{tabular}{lll}
\toprule
 & \file{ckpt.pt} (fat) & \file{model.pt} (slim) \\
\midrule
contents & model, optimizer, data cursor, & model, config, step, tokens\_seen \\
         & step, tokens\_seen, config & \\
size (base) & 1.48GB & 495MB \\
purpose & resume a training run & inference and release \\
\bottomrule
\end{tabular}
\caption{The two checkpoint formats. The optimizer's AdamW state is two extra numbers per
parameter, which is why the fat checkpoint is roughly 3x the slim one.}
\end{table}

\begin{keyidea}{Both files carry \code{config}, which is what makes eval scripts work}
\code{torch.save(\{..., "config": model\_cfg.\_\_dict\_\_\})} means a checkpoint is
self-describing. \code{eval/\_shared.py}'s \code{load\_model} does
\code{ModelConfig(**ck["config"])} and rebuilds the exact architecture from the file
alone, with no need to know which preset or which config file produced it. This is why
\code{python -m eval.perplexity --ckpt runs/tiny/model.pt} just works.
\end{keyidea}

\begin{honest}{The docstring promises RNG restoration that the checkpoint does not contain}
\file{train.py}'s module docstring advertises: "full resume: model + optimizer +
scheduler + data cursor + \textbf{RNG}". The checkpoint's actual keys are
\code{model}, \code{optimizer}, \code{data}, \code{step}, \code{tokens\_seen},
\code{config}. \textbf{No RNG state is saved or restored.}

How much does it matter? Less than it sounds. The data order is governed by the seeded
shuffle and the cursor, both of which \emph{are} restored, and dropout is 0.0 in every
committed config, so the only consumer of the RNG in the training path is the
initialization that a resume overwrites anyway. So the resume is very likely
bit-reproducible in practice despite the missing state. The defect is the claim, not the
behaviour. The README's own wording ("optimizer + scheduler + data cursor") is accurate;
only the docstring overreaches. Fix the docstring.

Note also that "scheduler" is restored in a subtle sense: there is no
\code{torch.optim.lr\_scheduler} object at all. The LR is recomputed from \code{step} by
\code{cosine\_lr} every iteration, so restoring \code{step} restores the schedule exactly.
That is a better design than a stateful scheduler, and it is worth pointing out rather
than letting it look like an omission.
\end{honest}

\begin{honest}{A crash between checkpoint intervals loses up to 2000 steps}
\code{ckpt\_interval} is 2000 for \code{base}, which at 524,288 tokens per step is over
1B tokens, roughly 2.8 hours of A100 time at the measured throughput. A crash at step
3,999 would resume from step 2,000. Given the v2 run cost \$13.06 in total, the exposure
is a few dollars, so the setting is defensible. But the number was evidently chosen for
disk-write frequency rather than by reasoning about the loss window, and 500 would have
cost almost nothing extra.
\end{honest}

\subsection{The configs}

\begin{table}[H]
\centering
\small
\begin{tabular}{lrrrrrr}
\toprule
 & \textbf{block} & \textbf{batch} & \textbf{accum} & \textbf{tokens/step} & \textbf{max\_tokens} & \textbf{lr} \\
\midrule
tiny   &  256 & 16 &  2 &   8,192 & 4.096M & 6e-4 \\
small  & 1024 & 24 &  8 & 196,608 & 2B     & 6e-4 \\
base   & 1024 & 32 & 16 & 524,288 & 450M   & 3e-4 \\
medium & 1024 & 32 & 16 & 524,288 & 3.45B  & 2.5e-4 \\
\bottomrule
\end{tabular}
\caption{The four committed configs. Learning rate falls as model size rises, which is
the conventional relationship.}
\label{tab:configs}
\end{table}

\begin{honest}{\code{configs/base.yaml} does not reproduce \code{runs/base\_v2}, and nothing in the repo records what would}
This is the most consequential documentation gap in the project.

The committed \file{configs/base.yaml} says \code{max\_tokens: 450000000}, and its
comment explains that as "\textasciitilde{}4 epochs of the 111.7M-token corpus". Both are v1-era facts.
The README states that \code{runs/base\_v2}, the project's headline model, trained on
\textbf{2,499,805,184} tokens. That is $4{,}768 \times 524{,}288$, confirmed exactly
against the committed \file{runs/base\_v2/metrics.jsonl}, whose last logged step is 4,760.

So the v2 run was launched with \code{--max-tokens} overriding the config, and the anneal
phase additionally needed \code{--data-dir data\_anneal} on a \code{--resume}. Both flags
exist in \code{train.py}. But \textbf{the exact commands are written down nowhere}:
\file{docs/GPU\_RUN.md} still documents only \code{python train.py --config
configs/base.yaml}, which would train a different model on a different budget.

Someone cloning this repository and following the documentation cannot reproduce the
model the repository is about. The fix is small: either commit a \file{base\_v2.yaml} and
an \file{anneal.yaml} with the real values, or write the two literal commands into
\file{GPU\_RUN.md}. The training code is not at fault; the record of how it was invoked
is missing.
\end{honest}

\begin{honest}{Two stale comments in the configs}
\file{configs/tiny.yaml}'s header says "\textasciitilde{}10M params"; the real figure is 13.01M, measured
above. \file{configs/base.yaml}'s \code{grad\_accum} comment says "effective \textasciitilde{}524k
tokens/step on a single 24GB GPU"; both real runs used an A100 SXM \textbf{80GB}. Neither
misleads about anything load-bearing, and the README is scrupulous about the 13M figure
elsewhere, so these are simply comments that were not revisited. \file{configs/medium.yaml},
by contrast, is exemplary: it explicitly flags its own \code{batch\_size} as carried over
from \code{base.yaml} and unverified against the larger model's memory footprint.
\end{honest}

\subsection{Distributed training}

\code{is\_dist()} checks the \code{WORLD\_SIZE} environment variable that
\code{torchrun} sets, so a plain \code{python train.py} takes the single-process path
with zero distributed overhead and the same file works under \code{torchrun} unchanged.
The backend is chosen by hardware: \code{nccl} on CUDA, \code{gloo} otherwise. Seeds are
offset per rank (\code{seed = cfg.get("seed", 1337) + rank}) so ranks do not draw
identical dropout masks.

\begin{jargon}{DDP (DistributedDataParallel)}
Each GPU holds a full copy of the model and processes a different slice of the batch.
After the backward pass the gradients are averaged across all GPUs, so every copy takes
an identical optimizer step and stays in sync. It scales throughput nearly linearly.
\end{jargon}

\begin{honest}{The DDP path was never exercised}
Both real runs used a single A100, and \code{docs/GPU\_RUN.md}'s log shows
\code{python3 train.py}, not \code{torchrun}. The DDP code is written carefully (the
gradient-sync suppression above is a detail that only someone who has thought about DDP
includes) but it is untested: no test covers it, and \code{world\_size > 1} has never run.
The \code{require\_backward\_grad\_sync} attribute in particular is set on \code{model},
which is the DDP wrapper in that branch, which is correct, but "correct by reading" is
not the same as "known to work". If asked, the honest answer is that the multi-GPU path
is written but unproven.
\end{honest}

\newpage
% =====================================================================
\section{Part VII: Evaluation}
\label{sec:eval}

Four scripts, measuring four different things. Together they answer: is it a good
language model (perplexity), is it writing Urdu at all (script mix), does it know
anything (cloze), and was the tokenizer worth building (compare, already covered in
Section~\ref{sec:tokclaim}).

\subsection{\code{eval/perplexity.py}: the primary metric}

Loads a checkpoint, memory-maps a \file{.bin} split, cuts it into blocks at fixed
offsets, and averages the loss.

\begin{lstlisting}[language=Python,caption={The token-weighted average}]
_, loss = model(x, y)
total_loss += loss.item() * x.numel()      # loss is a mean; re-weight by tokens
total_tokens += x.numel()
avg = total_loss / total_tokens
print(f"tokens={total_tokens} loss={avg:.4f} perplexity={math.exp(avg):.2f}")
\end{lstlisting}

\begin{keyidea}{Why the multiplication by \code{x.numel()} is correct and necessary}
\code{F.cross\_entropy} returns the \emph{mean} loss over the batch's tokens. Averaging
those means directly would be wrong if batches held different token counts (the last
batch usually does). Multiplying each mean back out by its token count and dividing by
the grand total gives a properly token-weighted average. It is a small thing that is
frequently got wrong.
\end{keyidea}

Unlike the loop's in-training eval, this script is fully deterministic: fixed starts, no
shuffle, no sampling. It must reproduce exactly, which is what made the check below
possible.

\begin{honest}{Verified, and it does not reproduce: 214.11, not the README's 248.8}
The README states the tiny proof run's "final held-out perplexity (fuller eval on
\file{data/val.bin}): \textbf{248.8}". Running the committed script against the
committed checkpoint and the committed data on 2026-07-16 gives:
\begin{lstlisting}
$ python -m eval.perplexity --ckpt runs/tiny/model.pt --bin data/val.bin
tokens=409600 loss=5.3665 perplexity=214.11

$ python -m eval.perplexity --ckpt runs/tiny/model.pt --bin data/val.bin \
      --max-batches 100000                      # the entire val split
tokens=1118208 loss=5.1762 perplexity=177.01
\end{lstlisting}
The script is deterministic, so this is not sampling noise. It was run twice, at the
default coverage and over the whole validation set, to rule out the possibility that
248.8 corresponded to some other batch budget. It does not: more coverage moves the
number \emph{further} from the claim, not closer. The 248.8 was produced from different
inputs.

The likely explanation is visible in the file timestamps. \file{runs/tiny/model.pt} was
written at 20:21, but \file{data/val.bin} and \file{data/corpus\_stats.json} were both
rewritten at \textbf{21:46}, over an hour \emph{later}. The validation set was
regenerated after the tiny model was trained and evaluated, so the committed
\file{val.bin} is not the one the 248.8 was measured on. That is also consistent with the
direction of the gap: the current number is \emph{better} than the claim, so nothing was
inflated.

This is worth fixing rather than shrugging at, for two reasons. It is the one number in
this repository found not to reproduce, in a project whose entire pitch is honest
reporting; and a reader who checks it will find a mismatch and start wondering about the
rest, which would be unfair to a project where everything else checked out to the digit.
Re-run the eval and update the README, or state which val.bin the figure came from.
\end{honest}

\begin{honest}{"Fuller eval" is actually the default, and it covers about 37\% of val}
The README calls the 248.8 a "fuller eval", implying more coverage than the in-training
one. But the command it gives passes no flags, so it uses the defaults
\code{--max-batches 200 --batch-size 8}. With the tiny model's 256-token blocks that is
$200 \times 8 \times 256 = 409{,}600$ tokens, confirmed by the script's own output above.
The committed \file{data/val.bin} is 2,236,612 bytes of uint16, so 1,118,306 tokens. The
"fuller" eval therefore covers about \textbf{37\%} of the validation set, and is capped
by \code{--max-batches}, not by the data. It is fuller than the in-training eval (which
uses \code{eval\_iters} batches) but it is not a full pass, and the phrase invites the
reader to think it is.
\end{honest}

\subsection{\code{eval/script\_mix.py}: is it even writing Urdu?}

A diagnostic no English-language project needs, and a genuinely clever cheap check.

\begin{lstlisting}[language=Python,caption={Classifying generated characters}]
for c in text:
    if not c.isalpha():   continue        # skip digits and punctuation
    if is_arabic_char(c): arabic += 1
    elif c.isascii():     latin += 1
    else:                 other += 1
\end{lstlisting}

It generates from five fixed one-word Urdu prompts, cycled by index, at a temperature of
0.9 with \code{top\_k} of 40, and reports the mean share of each script. It reuses
\code{is\_arabic\_char} from the pipeline's own \file{normalize.py}, so the diagnostic
and the corpus filter share exactly one definition of "Arabic script". That reuse is the
right instinct.

\begin{keyidea}{What this metric actually catches}
A model can have a plausible-looking perplexity while quietly degenerating into Latin
text it absorbed from crawl contamination, or into punctuation loops. Perplexity would
not obviously flag it; this would, instantly and in one number. The v2 model scores
99.75\% Arabic script over 50 generations, which says the corpus filtering worked and the
model learned to stay in Urdu.
\end{keyidea}

\begin{honest}{The prompt is included in the text being scored}
\code{text = tok.decode(out[0].tolist())} decodes the \emph{whole} sequence, prompt
included, and \code{script\_shares} then counts all of it. Since every prompt is Urdu, the
prompt's characters are counted as Arabic-script wins the model did not earn. With
\code{max\_new\_tokens=80} against a one-word prompt the contamination is a few percent
at most, so it cannot manufacture the 99.75\% result. But it does mean the metric is
biased upward by construction, and slicing off the prompt before scoring would be a
two-line fix. Worth knowing before someone asks whether the 99.75\% includes the prompt,
because it does.
\end{honest}

\subsection{\code{eval/cloze.py}: does it know anything?}
\label{sec:cloze}

\begin{jargon}{Cloze test}
A fill-in-the-blank. The model is given a stem (a sentence up to a blank) and several
candidate completions, exactly one correct. It does not generate; it \emph{scores} each
candidate and the best-scoring one is its answer. This measures knowledge without
requiring the model to be able to follow an instruction, which a 124M base model cannot.
\end{jargon}

\begin{jargon}{Few-shot prompting}
Prepending a few solved examples to the prompt so the model can infer the pattern from
context. \code{--shots 2} builds a prefix from the first two items (stem plus correct
answer) and prepends it to every subsequent question. Those two items are then excluded
from scoring, which is why 95 items yield 93 scored.
\end{jargon}

\begin{lstlisting}[language=Python,caption={Scoring one option}]
ctx_ids = tok.encode(context).ids
opt_ids = tok.encode(option).ids
ids = (ctx_ids + opt_ids)[-max_len:]
x = torch.tensor([ids[:-1]]); y = torch.tensor([ids[1:]])
mask_start = len(ids) - 1 - len(opt_ids)
y = y.clone()
y[0, :max(mask_start, 0)] = -100          # score ONLY the option's tokens
_, loss = model(x, y)
return loss.item()
\end{lstlisting}

The masking is the heart of it, and it is right: everything before the option is set to
\code{-100} so the context contributes nothing to the loss, and only the candidate's own
tokens are scored. \code{[-max\_len:]} crops from the left if the few-shot prefix has
overflowed the context window, keeping the option intact at the end where it must be.

\begin{honest}{The docstring says "total log-probability"; the code computes a length-normalized mean}
The module docstring says the harness scores "the total log-probability of the completion
tokens". It does not. \code{F.cross\_entropy}'s default reduction is \code{"mean"}, and
with \code{ignore\_index=-100} it averages over the non-ignored positions only. So
\code{option\_loss} returns the \textbf{mean per-token} loss of the option, not the sum.

The difference is not cosmetic. Total log-probability systematically favours \emph{short}
options, because every extra token adds more negative log-probability; mean per-token
loss removes that length bias. Since the eval set's distractors are described as
"grammatical and same-shape", the options are of similar length and the two would rank
them nearly identically here, so the reported numbers are almost certainly unaffected.

Which makes this a documentation bug rather than a scoring bug, and arguably the code
made the \emph{better} choice than the docstring describes. But it is precisely the kind
of detail an interviewer probes ("how do you handle length bias in multiple-choice
scoring?"), and the answer should be "mean per-token, deliberately", not a docstring that
says otherwise.
\end{honest}

\subsubsection{The eval set}

\begin{keyidea}{Verified: 95 items, 3 options each, hand-written}
\file{evaldata/cloze.jsonl} was read directly: exactly 95 lines, every item with exactly
3 options, fields \code{id}, \code{stem}, \code{options}, \code{answer\_idx},
\code{translation}. The README's "expanded from 45 to 95 items" is confirmed by the file.
\end{keyidea}

\begin{honest}{\file{evaldata/README.md} still says 45 items}
The eval set's own README describes "45 items" and "a probe for the phase-2 base model...
the tiny CPU proof-run checkpoint scores near chance". The file has held 95 items since
2026-07-08, and the base model has been trained twice since. The top-level README is
current; this one was not updated alongside it. It also still refers to the tiny run as
the reference point, which is two runs out of date.
\end{honest}

\begin{honest}{Every one of the 95 items has \code{answer\_idx} 0}
Verified by counting: the correct answer is the first option in all 95 items, with no
exceptions. \file{evaldata/README.md} discloses this ("placed first in every item for
authoring convenience") and correctly notes that the harness scores log-probabilities and
never looks at position, so there is no leak through the scoring path.

That defence is sound, and this is not a live bug. But it costs something real: it makes a
whole class of bug undetectable. If a refactor ever introduced a position bias, or if
\code{pred} were computed wrongly and happened to return 0, the harness would score 100\%
and look like a triumph. A shuffled \code{answer\_idx} would make that failure mode
impossible and costs one script to produce. For a set explicitly built to be trusted for
future comparisons, that is worth doing.
\end{honest}

\begin{honest}{The Wikidata pipeline would contaminate this eval set, and the code says otherwise}
This one matters, because it concerns a run that has not happened yet and can therefore
still be fixed.

\file{pipeline/fetch\_wikidata.py}'s docstring states the fact triples are "using
different entities than the eval set's own 45 items", and \file{docs/DATA\_LICENSES.md}
repeats it: "distinct entities from the eval set's own 45 cloze items". \textbf{Nothing
in the code enforces this, and it is demonstrably false.}

The \code{capital} query is \code{?subject wdt:P36 ?object . ?subject wdt:P31 wd:Q6256},
which selects the capital of \emph{every country}. Pakistan is a country. So the query
returns Pakistan/Islamabad, and \code{read\_wikidata\_facts} renders it into a
declarative Urdu sentence of the form "the capital of Pakistan is Islamabad".

Cloze item \code{id} 1, read directly from \file{evaldata/cloze.jsonl}, glosses in its own
\code{translation} field as: "The capital of Pakistan is Islamabad." Item 21 is "Urdu is
the national language of Pakistan", which the \code{official\_language} (P37) query
covers the same way.

So training on the Wikidata source as configured would put at least one eval item's exact
fact, in near-identical declarative phrasing, directly into the training data. That is
textbook train/test contamination, and it lands on the \emph{one metric} the \code{medium}
run exists to improve. A cloze gain measured that way would be partly unearned, and the
project's whole argument (that fact recall needs more parameters, per Roberts et al.) is
exactly the kind of claim that a contaminated eval would appear to confirm.

Two honest caveats, in the project's favour. First, the contamination is small: two of 95
items are affected by the four relations queried. Second, Urdu Wikipedia certainly already
contains prose stating Pakistan's capital, so the eval set was never contamination-free in
the strict sense, and no one claims a base model's cloze score is a clean held-out
measurement. What is specifically wrong is the \emph{documented claim of
entity-distinctness}, which is asserted twice and enforced nowhere.

The fix is easy and should happen before the \code{medium} run: load
\file{evaldata/cloze.jsonl}, extract its subjects, and filter them out of the fetched
triples, with a test pinning it. Then the docstring becomes true.
\end{honest}

\subsection{\code{eval/plot\_loss.py} and \code{sample.py}}

\code{plot\_loss.py} reads \file{metrics.jsonl} and renders the curve to a PNG with
matplotlib, using the \code{Agg} backend so it works on a headless box. It keys on the
presence of \code{"loss"} versus \code{"val\_loss"} to separate the two series, which is
exactly the shape the training loop writes.

\code{sample.py} loads a checkpoint and prints continuations for three default Urdu
prompts (transliterated: \emph{pakistan ek}, \emph{urdu zaban}, \emph{aaj mausam}).

\begin{honest}{\code{sample.py} duplicates \code{load\_model} verbatim from \code{eval/\_shared.py}}
The two functions are identical, line for line: same \code{torch.load}, same
\code{ModelConfig(**ck["config"])}, same \code{.eval()}, same return. \code{eval/cloze.py}
and \code{eval/script\_mix.py} both import the shared one; \code{sample.py} defines its
own copy instead. Presumably \file{sample.py} sits at the repo root rather than inside
\file{eval/} and the import was avoided. \code{conftest.py} already puts the repo root on
\code{sys.path}, and the eval scripts import \code{pipeline.common} across packages
anyway, so \code{from eval.\_shared import load\_model} would work. Fifteen duplicated
lines is trivial, but it is the one place where the repository's otherwise good factoring
slips.
\end{honest}

\newpage
% =====================================================================
\section{Part VIII: What actually happened when it ran}
\label{sec:runs}

Three real runs. Everything in this part was read out of the committed
\file{metrics.jsonl} files, not out of the README.

\subsection{The tiny CPU proof run}

The point was never capability. It was to prove the pipeline end to end on an 8-core CPU:
real corpus, real tokenizer, real training loop, a loss that actually falls, and a
checkpoint that loads.

\begin{figure}[H]
\centering
\begin{tikzpicture}
\begin{axis}[
  width=13.2cm, height=5.6cm,
  xlabel={step}, ylabel={cross-entropy loss (nats)},
  xmin=0, xmax=510, ymin=4.5, ymax=10,
  tick label style={font=\small}, label style={font=\small},
  xmajorgrids, ymajorgrids, grid style={linegrey},
  legend style={font=\footnotesize, at={(0.98,0.95)}, anchor=north east, draw=linegrey},
]
\addplot[draw=accent, thick, mark=none] coordinates {
(20,9.736) (40,8.639) (60,7.466) (80,6.991) (100,6.716) (120,6.494) (140,6.542)
(160,5.963) (180,6.144) (200,5.924) (220,5.806) (240,5.638) (260,5.528) (280,5.486)
(300,5.719) (320,5.418) (340,5.564) (360,5.378) (380,5.534) (400,4.964) (420,5.381)
(440,5.263) (460,5.169) (480,5.311) (500,5.420) };
\addplot[draw=warnred, thick, mark=*, mark size=2pt, only marks] coordinates {
(200,5.8678) (400,5.3572) };
\legend{train loss, held-out val loss}
\end{axis}
\end{tikzpicture}
\caption{The tiny run, drawn from the committed \file{runs/tiny/metrics.jsonl}. 500 steps,
4.10M tokens. The curve is still clearly falling at the end; the run was stopped at a
checkpoint boundary, not at convergence.}
\label{fig:tinycurve}
\end{figure}

Verified against the log: first logged train loss \textbf{9.7365} at step 20, last
\textbf{5.4204} at step 500; val loss \textbf{5.8678} at step 200 and \textbf{5.3572} at
step 400. These match the README's "9.74 to 5.42" and its val perplexities (352.5 and
212.1, which are $e^{5.8678}$ and $e^{5.3572}$) exactly.

Note the honesty the curve itself demands: the last three train points are 5.169, 5.311,
5.420, so the final logged value is a \emph{local uptick}, not the minimum. The README
quotes 5.42, the final value, rather than cherry-picking the 4.964 at step 400. That is
the correct choice and the harder one.

\subsection{Base v1: the first real GPU run}

\begin{table}[H]
\centering
\small
\begin{tabular}{ll}
\toprule
Date and hardware & 2026-07-07, rented A100 SXM 80GB (RunPod) \\
Steps and tokens & 858 steps, 449.8M tokens, roughly 4 epochs of the v1 corpus \\
Measured throughput & $\sim$99,470 tokens/second sustained \\
Wall clock and cost & about 75 minutes, about \$1.90 \\
Loss & first logged 9.5863 (step 20), last logged 3.5087 (step 840) \\
Held-out perplexity & 38.02 \\
\bottomrule
\end{tabular}
\caption{Base v1. Throughput and losses read from \file{runs/base/metrics.jsonl}.}
\end{table}

\begin{honest}{The README's "loss fell 7.35 to 3.51" understates the starting loss}
7.3513 is a real number in \file{runs/base/metrics.jsonl}, but it is the value at
\textbf{step 100}, not the start. The first logged loss is \textbf{9.5863} at step 20.
The run's honest range is 9.59 to 3.51.

Two things make this worth flagging despite being minor. It is inconsistent with the
project's own reporting of v2 in the same README paragraph, which correctly quotes 9.63
(the step-20 value) as that run's starting loss, so the two runs are described on
different bases and are not directly comparable as written. And it understates the
model's actual improvement, which is the opposite of the direction people usually err, so
this reads as an editing slip rather than anything else. The fix is one number.
\end{honest}

\begin{honest}{v1 has no in-training validation numbers at all}
\file{runs/base/metrics.jsonl} contains \textbf{zero} \code{val\_loss} records. The
reason is arithmetic: \code{eval\_interval} is 1000 in \file{configs/base.yaml}, and the
run's total was 858 steps, so the evaluation branch never fired once. The 38.02
perplexity in the README came from a separate \code{eval.perplexity} invocation
afterwards, which is fine and is what that script is for. But it means the v1 run flew
blind: there was no signal during those 75 minutes that would have revealed overfitting
or divergence. v2, at 4,768 steps, got four val checks. Setting \code{eval\_interval}
above the run's own step count is a small planning slip worth noticing.
\end{honest}

\subsection{Base v2: the headline model}

The v2 run is the project's real result: roughly 7.7x more unique tokens than v1, a
token/parameter ratio of about 20.2 against v1's 0.9, and a two-phase schedule ending in
a Wikipedia-only anneal.

\begin{figure}[H]
\centering
\begin{tikzpicture}
\begin{axis}[
  width=8.8cm, height=6.4cm,
  xlabel={step (524,288 tokens each)}, ylabel={train loss (nats)},
  xmin=0, xmax=4850, ymin=2.3, ymax=10,
  tick label style={font=\small}, label style={font=\small},
  xmajorgrids, ymajorgrids, grid style={linegrey},
  legend style={font=\footnotesize, at={(0.98,0.95)}, anchor=north east, draw=linegrey},
]
\addplot[draw=accent, thick, mark=none] coordinates {
(20,9.629) (140,6.831) (260,5.678) (380,5.027) (500,4.627) (620,4.235) (740,3.764)
(860,3.469) (980,3.174) (1100,3.200) (1220,3.154) (1340,3.052) (1460,3.000)
(1580,2.825) (1700,2.933) (1820,2.803) (1940,2.791) (2060,2.808) (2180,2.731)
(2300,2.725) (2420,2.686) (2540,2.675) (2660,2.718) (2780,2.704) (2900,2.625)
(3020,2.682) (3140,2.664) (3260,2.584) (3380,2.665) (3500,2.688) (3620,2.572)
(3740,2.583) (3860,2.605) (3980,2.584) (4100,2.496) (4220,2.612) (4340,2.550)
(4460,2.469) (4580,2.729) (4700,2.773) (4760,2.741) };
\draw[warnred, thick, dashed] (axis cs:4570,2.3) -- (axis cs:4570,10);
\node[lbl, anchor=south east, text=warnred, font=\scriptsize] at (axis cs:4540,6.4)
  {anneal starts};
\legend{train loss}
\end{axis}
%
\begin{axis}[
  at={(8.7cm,0)}, anchor=south west,
  width=6.0cm, height=6.4cm,
  xlabel={step (zoom)}, ylabel={},
  xmin=4390, xmax=4770, ymin=2.40, ymax=2.86,
  tick label style={font=\scriptsize}, label style={font=\small},
  xtick={4400,4500,4600,4700}, xticklabel style={/pgf/number format/fixed,
    /pgf/number format/1000 sep={,}},
  xmajorgrids, ymajorgrids, grid style={linegrey},
  title={\small the anneal switch, every logged point},
  title style={font=\small, text=inkblue},
]
\addplot[draw=accent, thick, mark=*, mark size=1.1pt] coordinates {
(4400,2.5372) (4420,2.5413) (4440,2.5392) (4460,2.4686) (4480,2.5343) (4500,2.5530)
(4520,2.5075) (4540,2.5336) (4560,2.5648) };
\addplot[draw=warnred, thick, mark=*, mark size=1.1pt] coordinates {
(4580,2.7288) (4600,2.7339) (4620,2.7435) (4640,2.8064) (4660,2.7618) (4680,2.7515)
(4700,2.7733) (4720,2.7443) (4740,2.6574) (4760,2.7414) };
\draw[warnred, thick, dashed] (axis cs:4570,2.40) -- (axis cs:4570,2.86);
\node[lbl, anchor=north east, text=warnred, font=\scriptsize] at (axis cs:4562,2.855)
  {Wikipedia-only\\ from here};
\node[lbl, anchor=south west, font=\scriptsize] at (axis cs:4398,2.415) {CC-100 mix};
\end{axis}

\end{tikzpicture}
\caption{The v2 run, every 6th logged point from the committed
\file{runs/base\_v2/metrics.jsonl}, plus the final point. The loss step at the anneal
boundary is real data, not an annotation.}
\label{fig:v2curve}
\end{figure}

\begin{keyidea}{Verified: the anneal is visible in the log, exactly where it should be}
The README says the loss "rose briefly and expectedly at the phase-1 to anneal
data-source switch, then settled". The committed log confirms this precisely. Steps 4400
to 4560 sit at 2.47 to 2.56; step 4580 jumps to 2.73 and the run finishes in the 2.74 to
2.81 band. The jump lands between steps 4560 and 4580.

That is exactly where arithmetic predicts it. The main phase was about 2.4B tokens, which
at 524,288 tokens per step is roughly step 4,578. So the log independently corroborates
the README's account of a two-phase run, and the rise is what a data-distribution change
looks like, not instability: Wikipedia-only text is a different (and for this model,
harder) distribution than the CC-100 mix it had adapted to.
\end{keyidea}

\begin{table}[H]
\centering
\small
\begin{tabular}{lrrrr}
\toprule
\textbf{step} & \textbf{tokens} & \textbf{val loss} & \textbf{val perplexity} \\
\midrule
1000 &   524,288,000 & 3.9621 & 52.57 \\
2000 & 1,048,576,000 & 3.5324 & 34.21 \\
3000 & 1,572,864,000 & 3.3818 & 29.42 \\
4000 & 2,097,152,000 & 3.3074 & 27.31 \\
\bottomrule
\end{tabular}
\caption{Every in-training validation record from \file{runs/base\_v2/metrics.jsonl},
matching the README's 52.57 and 27.31 endpoints. The curve is flattening: doubling the
tokens from 1B to 2B bought 6.9 perplexity points, while the next 1B bought 2.1.}
\label{tab:v2val}
\end{table}

\subsection{v1 against v2, and the honest reading}

\begin{table}[H]
\centering
\small
\begin{tabular}{lrr}
\toprule
\textbf{Metric} & \textbf{v1 (2026-07-07)} & \textbf{v2 (2026-07-08)} \\
\midrule
Unique training tokens & 111.7M & 862.4M \\
Tokens seen & 449.8M & 2,499,805,184 \\
Tokens per parameter & 0.9 & 20.2 \\
Held-out perplexity & 38.02 & \textbf{28.93} \\
Script mix (Arabic) & 99.43\% (20 gens) & \textbf{99.75\%} (50 gens) \\
Cloze, 45-item set & 34.9\% (15/43) & 39.5\% (17/43) \\
Cloze, 95-item set & not measured & 45.2\% (42/93) \\
GPU time and cost & $\sim$75 min, $\sim$\$1.90 & 6h 39m, $\sim$\$13.06 \\
\bottomrule
\end{tabular}
\caption{The two GPU runs. Cloze figures are README-reported; the perplexity, script-mix
and cloze numbers for the base models could not be re-run here (see below).}
\label{tab:v1v2}
\end{table}

\begin{keyidea}{The project's own reading of these numbers is the correct one}
The README does something unusual and right: it refuses to claim the cloze improvement.
34.9\% to 39.5\% is a 2-item difference on a 43-item set whose standard error is roughly
$\pm$7 points, so it sits well inside noise. The README says so in as many words, rather
than presenting "+4.6 points" as a result.

It then explains \emph{why} cloze barely moved while perplexity fell 24\%, and the
explanation is both correct and well-sourced: a 124M model has a real ceiling on how many
discrete facts it can store, and that capacity scales with parameter count, not with how
much text it reads (Roberts et al., 2020). So 7.7x more text made it more \emph{fluent}
(which is what perplexity measures) without making it know more (which is what cloze
measures). The corpus expansion did not fail; it did exactly what more text does.

That diagnosis then drives the plan: the \code{medium} preset exists because more
parameters is the lever that moves fact storage, and the eval set was widened from 45 to
95 items to cut the noise band from $\pm$7 to $\pm$5.2 precisely so the next comparison
can be trusted. Identifying that your headline metric is too noisy to support your next
claim, and fixing the metric before running the experiment, is the most mature single
decision in this repository.
\end{keyidea}

\begin{honest}{What could not be verified here, and why}
The base models' own evaluation numbers (perplexity 38.02 and 28.93, script mix, cloze)
were \textbf{not} re-run for this document, and are reported on the README's authority.
Two hard blockers:
\begin{itemize}
\item \file{data/train.bin} is gitignored (214MB for v1, 1.7GB for v2, both over
      GitHub's 100MB per-file limit), and the v2 corpus cannot be rebuilt without
      re-downloading roughly 1.3GB of raw sources and running a multi-hour pipeline.
\item The \code{base} weights are present via Git LFS, but evaluating a 124M model over
      the validation set on this 8-core CPU is hours of work, and the perplexity figures
      were measured against the \emph{v2} validation set, which is not the committed
      \file{data/val.bin} (that one is v1's, as the \file{corpus\_stats.json} shows).
\end{itemize}
What \emph{was} verified from the committed artifacts: the training losses, throughputs,
step and token counts, the anneal boundary, all four presets' parameter counts, the full
test suite, the eval set's shape, and the tokenizer's headline number. The base models'
downstream eval metrics remain README-attributed, and are labelled as such in
Table~\ref{tab:v1v2}.
\end{honest}

\begin{honest}{\file{docs/GPU\_RUN.md} was never updated after either run and is now actively misleading}
The document opens with "phase 2 is a budget decision, not an engineering one", estimates
throughput at "roughly 25k-40k tokens/s... taking a conservative 30k tok/s", and states
flatly: \emph{"These are estimates, not measured; the only measured throughput in this
repo is the CPU proof run."}

Two GPU runs have since happened, and the committed logs show a real sustained
\textbf{99,470} tok/s (v1) and \textbf{105,000} tok/s (v2). The "conservative" estimate
was low by a factor of about 3.3. Every downstream figure in that document inherits the
error: it projects the 450M-token run at $\sim$4.2 hours for about \$6, where v1 actually did
449.8M tokens in about 75 minutes for about \$1.90.

Its token-budget section is v1-era throughout ("our cleaned corpus is 111.7M training
tokens... only about 0.9 tokens per parameter"), which the v2 corpus superseded by 7.7x,
and its "Add more data before spending GPU money" advice was taken and completed on
2026-07-08. It even dates its own run log 2026-07-06 where the README says the v1 run was
2026-07-07.

The parts of that document written \emph{after} the run are excellent and are the reason
this is worth fixing rather than deleting: the RunPod SSH key injection trap, the warning
that installing the pinned \code{torch==2.12.1} would silently replace the pod's CUDA
build with a CPU-only one, the \code{tmux} detach pattern, and the note that SSH cannot
stop the billing meter so you must Terminate (not Stop) from the dashboard. That is
hard-won operational knowledge that no tutorial contains.

The document is simply two documents stapled together: a pre-run plan whose predictions
were falsified, and a post-run operations log that is entirely accurate. Splitting them,
and replacing the estimates with the measured numbers, would turn the project's weakest
document into one of its strongest. As it stands, a reader who trusts it will believe the
model is data-limited at 111.7M tokens and that an A100 does 30k tok/s, and both were
disproven by this project's own runs.
\end{honest}

\newpage
% =====================================================================
\section{Part IX: The planned \code{medium} run}

Three pieces are built and tested but have never been run: the \code{medium} preset
(336.38M parameters, verified above), the Wikidata fact pipeline, and the widened
95-item eval set. The reasoning chain behind them is the strongest argument in the
project, and it is worth being able to state it in one breath:

\begin{keyidea}{The argument, end to end}
v2's corpus expansion cut perplexity 24\% but left cloze inside the noise band. Roberts
et al. (2020) found fact-storage capacity scales with \emph{parameter count}, not with
training-text volume. Therefore more text cannot fix fact recall, and the next lever must
be a bigger model plus training data shaped as dense declarative facts. Hence
\code{medium} (2.7x the parameters) and the Wikidata triple pipeline. And because the old
45-item eval could not resolve the difference anyway, widen it to 95 first.
\end{keyidea}

The cost estimate in the README is explicitly labelled a projection, and its uncertainties
are named rather than buried: throughput is scaled from v2's measured $\sim$105,000 tok/s by
the parameter ratio ($336/124 \approx 2.71$, projecting $\sim$38,700 tok/s), which assumes
compute-bound scaling; and \file{configs/medium.yaml}'s \code{batch\_size: 32} is
inherited from \code{base.yaml} and unverified against a model whose activations are
roughly 2.3x larger, where \code{base} already used 36GB of 80GB. The recommendation to
smoke-test a few hundred steps before committing the $\sim$\$40-45 budget is the right one.

\begin{honest}{The one thing to fix before launching}
The Wikidata contamination described in Section~\ref{sec:cloze}. The \code{medium} run
exists to move the cloze number; feeding it training data containing at least two of the
eval set's own facts, while the code's own docstring claims the entities are distinct,
would compromise the exact measurement the run is for. It is perhaps thirty lines of
filtering and a test.
\end{honest}

\newpage
% =====================================================================
\section{Part X: Honest findings, collected}
\label{sec:honest}

Every finding in this document, gathered in one place and ordered by how much it would
cost to be surprised by it. Nothing here is fatal; this is a well-built repository and
the headline claims that could be checked all checked out. But an interviewer will find
some of these, and it is better to get there first.

\subsection{Substantive: worth fixing before the next run}

\begin{enumerate}
\item \textbf{Upsampling repeats documents back-to-back inside a single context window}
      (Section~\ref{sec:tokenizer}). \code{buf.extend(ids * reps)} lays three copies of a
      Wikipedia paragraph down contiguously, and at $\sim$70 tokens per paragraph all three
      fit inside one 1024-token block. This teaches copying rather than reweighting the
      distribution, and it shipped in the run that produced the headline model. Cheap to
      fix at the document-list level.
\item \textbf{The Wikidata fact pipeline would contaminate the cloze eval set, and two
      documents claim otherwise} (Section~\ref{sec:cloze}). The P36 query returns every
      country's capital, including Pakistan's, which is cloze item 1 verbatim. Nothing
      enforces the documented entity-distinctness. Must be fixed before the
      \code{medium} run, whose entire purpose is to move that metric.
\item \textbf{\code{configs/base.yaml} cannot reproduce \code{runs/base\_v2}, and nothing
      records what would} (Section~\ref{sec:training}). The config says 450M tokens; the
      run did 2.5B plus an anneal on a different data directory. The flags exist; the
      invocations are written down nowhere.
\item \textbf{The tiny run's reported perplexity does not reproduce} (Section~\ref{sec:eval}).
      README says 248.8; the committed script on the committed checkpoint and data gives
      \textbf{214.11} at the default coverage and \textbf{177.01} over the whole val
      split, so no batch budget produces the claim. The timestamps suggest
      \file{val.bin} was regenerated an hour after the figure was measured. The only
      number in the project found not to reproduce.
\item \textbf{\file{docs/GPU\_RUN.md} is a pre-run plan whose predictions were falsified,
      never updated} (Section~\ref{sec:runs}). It estimates 30k tok/s where the real runs
      measured 99.5k and 105k, states no GPU throughput has been measured when two runs
      have, and reasons throughout from the superseded 111.7M-token v1 corpus. Its
      post-run operations notes, by contrast, are excellent.
\end{enumerate}

\subsection{Real but contained}

\begin{enumerate}\setcounter{enumi}{5}
\item \textbf{The LSH index and the exact-hash set are held entirely in RAM}
      (Section~\ref{sec:dedup}), the one place the streaming discipline breaks. The
      README names this itself.
\item \textbf{\code{generate} has no KV cache}, so sampling is quadratic
      (Section~\ref{sec:model}). Fine for this project's usage, disqualifying for
      serving, and the reason the GQA branch has nothing to optimize.
\item \textbf{GQA is implemented with \code{repeat\_interleave}}, which materializes the
      expanded tensors and therefore saves no memory. It is a correctness shim, not an
      optimization, and no preset enables it. The README's word "capable" is exactly
      right.
\item \textbf{The DDP path has never been executed} (Section~\ref{sec:training}). Written
      carefully, including the gradient-sync suppression, but unproven and untested.
\item \textbf{The epoch wrap does not reshuffle and block offsets are fixed}, so v2's
      roughly 3 passes saw an identical block sequence each time.
\item \textbf{Stage resume is coarse}: markers are written only on stage completion, so a
      crash 90\% through extraction redoes all of it.
\item \textbf{v1 never ran a single in-training validation}, because
      \code{eval\_interval} (1000) exceeded the run's total steps (858).
\item \textbf{\code{script\_mix} scores the prompt along with the generation}, biasing the
      Arabic share upward. Too small to explain the 99.75\%, but real.
\item \textbf{The tokenizer comparison samples val-then-train}, and is held-out only
      because val happens to be big enough, not because anything enforces it.
\end{enumerate}

\subsection{Documentation and hygiene}

\begin{enumerate}\setcounter{enumi}{14}
\item \textbf{The README's tokenizer command omits \code{--max-chars 800000}} and as
      written yields 5.14x, not the 5.23x printed beside it (Section~\ref{sec:tokclaim}).
      The number is genuine and reproduces exactly with the flag.
\item \textbf{\code{cloze.py}'s docstring says "total log-probability"} but the code
      computes a length-normalized mean, which is the better choice, undocumented
      (Section~\ref{sec:cloze}).
\item \textbf{\code{train.py}'s docstring promises RNG restoration} that the checkpoint
      does not contain. Behaviourally near-harmless; the claim is the defect.
\item \textbf{README's v1 "loss fell 7.35"} is the step-100 value, not the step-20 start
      of 9.5863, and is inconsistent with how v2 is reported in the same paragraph.
\item \textbf{\file{evaldata/README.md} still says 45 items}; the file has held 95 since
      2026-07-08. \file{DATA\_LICENSES.md} likewise says "the eval set's own 45 cloze
      items".
\item \textbf{"Fuller eval" covers about 37\% of val}, capped by \code{--max-batches},
      not by the data.
\item \textbf{All 95 cloze items have \code{answer\_idx} 0}. Disclosed, and the harness is
      position-blind, but it makes a whole class of bug undetectable.
\item \textbf{\code{num\_params(non\_embedding=True)} has two identical branches}
      (Section~\ref{sec:model}); collapses to one line.
\item \textbf{\code{sample.py} duplicates \code{load\_model}} verbatim from
      \code{eval/\_shared.py}.
\item \textbf{\code{build\_minhash}'s \code{num\_perm} parameter is a trap}: the module's
      \code{\_A}/\code{\_B} are fixed at 64, so any other value would silently mismatch.
\item \textbf{\code{<|pad|>} and \code{<|unk|>} are structurally unreachable}
      (Section~\ref{sec:tokenizer}). Two inert vocabulary slots.
\item \textbf{The RoPE cache keys on \code{T} but ignores it}, so \code{generate} builds
      one identical full-size cache entry per token sampled.
\item \textbf{Stale config comments}: \file{tiny.yaml} says "\textasciitilde{}10M params" (real: 13.01M);
      \file{base.yaml} says "single 24GB GPU" (real: A100 80GB).
\item \textbf{\file{data/corpus\_stats.json} describes v1 only}; there is no committed
      stats file for the v2 corpus the headline model trained on.
\item \textbf{Pinned versions do not match the installed environment}
      (Section~\ref{sec:repro}). All 13 tests pass anyway.
\end{enumerate}

\subsection{What is genuinely strong}

It would be a distorted document that listed thirty findings and no strengths, so:

\begin{itemize}
\item \textbf{The headline claim is real and reproduces to the digit.} 806,495 chars,
      197,758 tokens, 4.078 chars/token, GPT-2's 1,034,831 and 0.779, gain 5.23x. Every
      figure, exactly.
\item \textbf{The MinHash vectorization} is a real 10x on a real bottleneck, with a
      correct and articulable reason why swapping the hash function is safe.
\item \textbf{The refusal to claim the cloze improvement}, followed by widening the eval
      set to fix the noise band \emph{before} running the experiment that needs it.
\item \textbf{The refusal to shrink the vocabulary} to make "10M params" true.
\item \textbf{The 8/3 SwiGLU ratio and the 14-layer consequence}, worked out rather than
      copied.
\item \textbf{Token-budgeted training}, which makes the Chinchilla reasoning directly
      expressible in the config and comparable across hardware.
\item \textbf{The causal-masking test}, which tests the property by experiment and
      includes the negative assertion that stops it being a tautology.
\item \textbf{The ZWNJ decision and the harakat decision}, both of which require actually
      knowing something about Urdu orthography rather than about machine learning.
\item \textbf{Validation is never upsampled or filtered}, and the tokenizer is trained on
      train only. Both leaks avoided quietly.
\item \textbf{The committed \file{metrics.jsonl} files}, which are what made most of this
      document checkable rather than merely repeatable.
\end{itemize}

\newpage
% =====================================================================
\section{Part XI: Reproduction, and what was verified}
\label{sec:repro}

\subsection{Dependency roles}

\begin{table}[H]
\centering
\small
\begin{tabular}{lll}
\toprule
\textbf{Package} & \textbf{Pinned} & \textbf{Role} \\
\midrule
torch            & 2.12.1  & the tensors, autograd, AdamW, SDPA/Flash, DDP \\
tokenizers       & 0.23.1  & Hugging Face BPE trainer and encoder (Rust) \\
numpy            & 2.5.1   & memmap, the vectorized MinHash, batch assembly \\
mwparserfromhell & 0.7.2   & MediaWiki wikitext parsing \\
datasketch       & 2.0.0   & MinHash LSH index \\
xxhash           & 3.8.1   & fast deterministic shingle hashing \\
PyYAML           & 6.0.3   & config loading \\
matplotlib       & 3.11.0  & the loss-curve PNG \\
pytest           & 9.1.1   & the 13 tests \\
tqdm             & 4.68.3  & \textbf{declared but never imported anywhere} \\
\bottomrule
\end{tabular}
\caption{\file{requirements.txt}. Note the last row.}
\label{tab:deps}
\end{table}

\begin{keyidea}{The dependency list is notably restrained}
There is no \code{transformers}, no \code{datasets}, no \code{accelerate}, no
\code{lightning}, no \code{wandb}. The modeling is hand-written, the data pipeline is
hand-written, the training loop is hand-written, and metrics go to a JSONL file. The only
heavyweight import is PyTorch itself. For a project whose purpose is to demonstrate that
the author can build this stack rather than assemble it, that restraint \emph{is} the
deliverable.
\end{keyidea}

\begin{honest}{\code{tqdm} is declared and never used}
\code{grep} finds no import of it anywhere in the repository. It is one line in
\file{requirements.txt} and harmless, but a reviewer reading the dependency list will
form an expectation of progress bars that the code does not meet.
\end{honest}

\begin{honest}{The pinned versions are not the versions this was verified against}
\file{requirements.txt} pins \code{torch==2.12.1}, \code{tokenizers==0.23.1},
\code{numpy==2.5.1}. The environment on this machine has torch 2.10.0+cu128,
tokenizers 0.21.1, numpy 1.26.4, and \textbf{all 13 tests pass and every measurement in
this document reproduced} regardless. So the code is more version-tolerant than the pins
suggest. \file{docs/GPU\_RUN.md} already documents the one place this bites: on a RunPod
pod, installing the pinned torch would silently replace the working CUDA build with a
CPU-only wheel, and the workaround is to filter torch out of the requirements before
installing. Exact pins plus a documented "except do not install this one" is a slightly
awkward combination; the pins are doing less work than they appear to.
\end{honest}

\subsection{What was verified for this document, and how}

\begin{table}[H]
\centering
\small
\begin{tabular}{p{6.6cm}p{2.3cm}p{5.2cm}}
\toprule
\textbf{Claim} & \textbf{Status} & \textbf{Evidence} \\
\midrule
Tokenizer: 5.23x vs GPT-2, 4.078 vs 0.779 chars/token, on an 806,495-char sample
  & \textbf{verified} & ran \code{eval.tokenizer\_compare} against a cached real GPT-2
    tokenizer, \code{--max-chars 800000}; every digit matched \\
13 tests pass (5 model, 8 pipeline) & \textbf{verified} & \code{pytest tests/ -q}:
    13 passed in 55.42s \\
Preset parameter counts 13.0 / 35.7 / 123.7 / 336.38M & \textbf{verified} &
    instantiated all four at vocab 32000 and counted \\
v1 corpus funnel (3.7M paragraphs, 31.94\% dedup, 2,489,976 / 25,106 split) &
    \textbf{verified} & read \file{data/corpus\_stats.json} \\
Tiny run: loss 9.74 to 5.42, val 5.87 to 5.36 & \textbf{verified} &
    read \file{runs/tiny/metrics.jsonl} \\
v1: 858 steps, 449.8M tokens, $\sim$99,470 tok/s & \textbf{verified} &
    read \file{runs/base/metrics.jsonl} \\
v2: 4,768 steps, 2,499,805,184 tokens, $\sim$105,000 tok/s, loss to $\sim$2.74 &
    \textbf{verified} & read \file{runs/base\_v2/metrics.jsonl}; token count is exactly
    $4768 \times 524288$ \\
v2 anneal caused a brief loss rise & \textbf{verified} & visible at step 4560 to 4580
    in the log, where arithmetic predicts it \\
v2 val perplexity 52.57 to 27.31 & \textbf{verified} & the four \code{val\_loss}
    records in the v2 log \\
Cloze set is 95 items, 3 options & \textbf{verified} & read
    \file{evaldata/cloze.jsonl} \\
\midrule
Tiny final perplexity 248.8 & \textbf{contradicted} & re-ran the script two ways:
    214.11 at the default coverage, 177.01 over the whole val split. Neither is 248.8 \\
\midrule
v1/v2 base perplexity (38.02 / 28.93), script mix, cloze accuracy &
    README-attributed & \file{train.bin} gitignored, v2 val set not committed, CPU
    infeasible \\
v2 corpus table (11.0M CC-100 docs, 37.5\% dedup, 13.3M train docs) &
    README-attributed & no committed stats file for v2 \\
GPU costs (\$1.90, \$13.06) and wall-clock & README-attributed & billing is external to
    the repository \\
890 Wikidata triples fetched, 865 surviving dedup & README-attributed & requires a live
    SPARQL call \\
\bottomrule
\end{tabular}
\caption{The verification log for this document. "Verified" means the code was run or the
committed artifact was read on 2026-07-16.}
\label{tab:verif}
\end{table}

\subsection{Running it yourself}

\begin{lstlisting}[caption={The paths that were exercised while writing this document}]
# tests
python -m pytest tests/ -q

# the headline claim (note the flag the README omits)
python -m eval.tokenizer_compare \
    --gpt2 ~/.cache/huggingface/hub/models--gpt2/snapshots/*/tokenizer.json \
    --max-chars 800000

# the tiny model, on the committed checkpoint and val split
python -m eval.perplexity --ckpt runs/tiny/model.pt --bin data/val.bin
python sample.py --ckpt runs/tiny/model.pt --seed 7
\end{lstlisting}

Everything above works from a fresh clone with \code{pip install -r requirements.txt},
because \file{data/val.bin}, \file{tokenizer/tokenizer.json} and
\file{runs/tiny/model.pt} are all committed. Rebuilding the corpus requires roughly 1.3GB
of raw downloads and several hours; the \code{base} weights need \code{git lfs pull}.

\begin{keyidea}{The last word}
The claim this project is built on, that an Urdu-native tokenizer is worth building
because English tokenizers shred the script, was checked here against a real GPT-2
tokenizer and reproduced exactly: 5.23x. The engineering underneath it is sound, the
scaling reasoning is correct and well-sourced, and the reporting is unusually honest,
including where it is unflattering.

The findings in Part X are almost entirely about \emph{documentation drifting behind
code}: a GPU-run plan never updated after the run, configs that no longer describe the
model that shipped, docstrings promising slightly more than they deliver, and a
perplexity measured against a validation set that was later regenerated. The code is in
better shape than the prose around it, which is the opposite of the usual failure and is a
much easier problem to fix.
\end{keyidea}

\end{document}
